% ============================================================================
% SETA-HPD, explained from first principles.
%
% This is a TUTORIAL document, not a conference paper. Its companion,
% sta_hpd_paper.tex, is written for a reader who already knows what split
% conformal prediction, exchangeability, and a highest-density region are; it
% is dense on purpose, because that is what a conference reviewer wants. This
% document assumes none of that. It starts from "what is a prediction
% interval and why would you want one," builds every concept in order
% (Gaussian density -> mixture-of-experts -> conformal prediction ->
% highest-density regions -> STA-HPD -> SETA-HPD), and only ever uses a
% definition, theorem, or piece of notation after it has been introduced.
% Every abstract construction is paired with a single running worked numeric
% example (a toy 2-expert model) so a reader can check every formula by hand.
% ============================================================================
\documentclass[11pt]{article}

\usepackage[margin=1in]{geometry}
\usepackage{amsmath, amssymb, amsfonts, amsthm}
\usepackage{bm}
\usepackage{booktabs}
\usepackage{algorithm}
\usepackage{algpseudocode}
\usepackage{enumitem}
\usepackage{hyperref}

\newtheorem{definition}{Definition}[section]
\newtheorem{example}{Example}[section]
\newtheorem{theorem}{Theorem}[section]
\newtheorem{lemma}{Lemma}[section]
\newtheorem{proposition}{Proposition}[section]
\newtheorem{corollary}{Corollary}[section]
\newtheorem{remark}{Remark}[section]

\title{\Large SETA-HPD, Explained from First Principles:\\
A Tutorial on Highest-Density-Region Conformal Calibration\\
for Mixture-of-Experts Regression}
\author{Gilad Aviv\\ \texttt{giladaviv987@gmail.com}}
\date{}

\begin{document}
\maketitle

\begin{abstract}
This document explains SETA-HPD -- a method for turning a Mixture-of-Experts regression model's
raw predictions into statistically valid, tightly-fitted prediction intervals -- starting from
nothing. It assumes no prior exposure to conformal prediction, highest-density regions, or
mixture models. Every idea is introduced once, defined precisely, motivated with plain-language
intuition, and then carried through a single running numeric example so that every formula in
the document can be checked by hand. The document ends exactly where the companion conference
paper (\texttt{sta\_hpd\_paper.tex}) begins: at a full statement of SETA-HPD, its two ancestor
methods (MoG-HPD and STA-HPD) as special cases, and the empirical results that motivate it. This
is a tutorial, not a novel contribution in itself -- everything here is a fuller explanation of
material that already exists in this repository's calibrators
(\texttt{calibration/tabular\_regression/MoG\_HPD\_Calibrator.py},
\texttt{STA\_HPD\_Calibrator.py}, \texttt{SETA\_HPD\_Calibrator.py},
\texttt{MELD\_HPD\_Calibrator.py}) and its conference-style write-up.
\end{abstract}

\tableofcontents

% ============================================================================
\section{Why do we need any of this?}
\label{sec:why}

\subsection{Point predictions are not enough}

Suppose a trained regression model looks at an input $\bm x$ (some features -- sensor readings,
past values of a time series, demographic variables, whatever the problem calls for) and
outputs a single number $\hat y(\bm x)$, its best guess at the true outcome $y$. This is called a
\emph{point prediction}. Point predictions are useful, but they are silent about a question that
matters enormously in practice: \emph{how wrong might $\hat y(\bm x)$ be?} A model that predicts
tomorrow's temperature as $21^\circ$C is not equally useful in a case where the true temperature
is almost always within $1^\circ$C of that guess, versus a case where it could plausibly be
anywhere from $10^\circ$C to $32^\circ$C. The point prediction looks identical in both cases; the
right decision to make downstream (should you bring a coat? should a hospital keep extra beds on
standby?) is completely different.

What we want instead of a single number is a \emph{prediction set} $\hat C(\bm x) \subseteq
\mathbb R$: a set of plausible values for $y$, built so that it is likely to actually contain the
true outcome. The simplest kind of prediction set is a symmetric interval $\hat y(\bm x) \pm
s(\bm x)$, but as this document will show in detail (Section~\ref{sec:hpd}), a symmetric interval
is not always the best-shaped set -- one of the two central ideas this document builds up to.

\subsection{What ``a good prediction set'' means}
\label{sec:coverage}

Before building any specific method, we need to pin down precisely what makes a prediction set
\emph{good}. Two properties, in tension with each other, define the whole subject:

\begin{itemize}[leftmargin=1.5em]
\item \textbf{Validity (coverage).} Fix a small number $\alpha \in (0,1)$, e.g.\ $\alpha=0.1$. A
prediction set $\hat C(\bm x)$ has \emph{marginal coverage} $1-\alpha$ if, averaged over the
randomness of drawing a fresh test point $(\bm x_{\mathrm{test}}, y_{\mathrm{test}})$,
\begin{equation}
P\big(y_{\mathrm{test}} \in \hat C(\bm x_{\mathrm{test}})\big) \ \ge\ 1-\alpha.
\label{eq:coverage-def}
\end{equation}
For $\alpha=0.1$ this says: if you built this prediction set for many independent test points,
the true outcome would fall inside the set at least $90\%$ of the time.
\item \textbf{Efficiency (width).} Among all sets with coverage at least $1-\alpha$, we prefer the
one with the smallest size (for an interval, the smallest width; more generally, the smallest
Lebesgue measure -- ``total length,'' allowing for a set made of several disjoint pieces).
\end{itemize}

These two properties trade off against each other in the most trivial way possible: the whole
real line $\hat C(\bm x) = \mathbb R$ always has coverage $1$, but it is useless (infinite width).
The entire technical content of every method in this document is about achieving \emph{valid}
coverage (Eq.~\ref{eq:coverage-def}) with the \emph{smallest possible width}. Section~\ref{sec:cp}
builds the tool that guarantees validity (conformal prediction); Section~\ref{sec:hpd} and
everything after it is about squeezing out width without ever giving up that validity guarantee.

\subsection{The model we are calibrating: a Mixture-of-Experts regressor}
\label{sec:moe}

Everything in this document is built on top of an already-trained \emph{Mixture-of-Experts} (MoE)
regression model, so we describe that model first, from scratch, before touching calibration at
all.

\subsubsection{A single Gaussian expert}

The simplest ``probabilistic'' regressor does not just output a number $\hat y(\bm x)$; it
outputs a full Gaussian (normal) density over possible outcomes,
\begin{equation}
p(y \mid \bm x) = \mathcal N(y;\ \mu(\bm x),\ \sigma^2(\bm x)) = \frac{1}{\sigma(\bm x)\sqrt{2\pi}}
\exp\!\left(-\frac{(y-\mu(\bm x))^2}{2\sigma^2(\bm x)}\right),
\end{equation}
where $\mu(\bm x)$ is the model's best-guess mean (its point prediction) and $\sigma(\bm x)$ is
its reported standard deviation -- how spread out it believes the outcome could be around that
mean, for this specific input $\bm x$. Because $\sigma(\bm x)$ is allowed to depend on $\bm x$,
this kind of model can be \emph{heteroscedastic}: more confident (small $\sigma$) on some inputs,
less confident (large $\sigma$) on others.

\subsubsection{Mixing several experts}

A single Gaussian is unimodal by construction -- it always has exactly one peak. Many real
outcomes are not: a delivery time might genuinely be ``fast, if there's no traffic'' or ``slow, if
there is,'' with little mass in between. A \emph{Mixture-of-Experts} model \cite{jordan1994hierarchical} represents this by
training $K$ separate sub-networks (``experts''), each of which looks at the same input $\bm x$
and outputs its own Gaussian $\mathcal N(\mu_k(\bm x), \sigma_k^2(\bm x))$, plus one more
sub-network (the \emph{gate}) that outputs, for the same input, a probability weight $w_k(\bm x)$
for each expert, with $\sum_{k=1}^K w_k(\bm x) = 1$. The model's full predictive density is the
weighted sum
\begin{equation}
f_{\mathrm{mix}}(y \mid \bm x) = \sum_{k=1}^K w_k(\bm x)\, \mathcal N\big(y;\ \mu_k(\bm x),\
\sigma_k^2(\bm x)\big).
\label{eq:mixdef}
\end{equation}
Read this as: ``with probability $w_k(\bm x)$, the outcome behaves like expert $k$'s Gaussian.''
Equation~\eqref{eq:mixdef} is a genuine probability density (it integrates to $1$, since each
component does and the weights sum to $1$), and unlike a single Gaussian it can be multimodal --
it can have several separate peaks, one per well-separated expert.

The model's single point prediction (used, e.g., to report ``the'' forecast, or to compute a mean
squared error) is the gate-weighted average of the expert means,
\begin{equation}
\hat y(\bm x) = \sum_{k=1}^K w_k(\bm x)\, \mu_k(\bm x).
\label{eq:yhatdef}
\end{equation}

\subsubsection{Two kinds of uncertainty: aleatoric and epistemic}
\label{sec:aleaepis}

A basic and useful fact from probability (the \emph{law of total variance}) splits the mixture's
variance around $\hat y(\bm x)$ into two pieces that have a genuine, different practical meaning:
\begin{align}
\underbrace{\operatorname{Var}(Y\mid\bm x)}_{\text{total variance}}
&= \underbrace{\sum_{k=1}^K w_k(\bm x)\,\sigma_k^2(\bm x)}_{\sigma_a^2(\bm x):\ \text{aleatoric (within-expert)}}
\ +\ \underbrace{\sum_{k=1}^K w_k(\bm x)\,\big(\mu_k(\bm x) - \hat y(\bm x)\big)^2}_{\sigma_e^2(\bm x):\ \text{epistemic (between-expert)}}.
\label{eq:lotv}
\end{align}
\begin{itemize}[leftmargin=1.5em]
\item $\sigma_a^2(\bm x)$, the \textbf{aleatoric} term, is how spread-out each individual expert
believes the outcome is, averaged over the gate's weighting. This is irreducible noise \emph{even
if the gate committed fully to one expert}.
\item $\sigma_e^2(\bm x)$, the \textbf{epistemic} term, is how much the experts \emph{disagree
with each other} -- how far each expert's mean sits from the aggregate prediction, weighted by how
much the gate trusts that expert. This is zero only when every expert the gate puts weight on
predicts (almost) the same mean.
\end{itemize}
We write $\sigma_{\mathrm{tot}}^2(\bm x) = \sigma_a^2(\bm x) + \sigma_e^2(\bm x)$ for the sum.

\begin{example}[The running example]
\label{ex:running}
Throughout this document we use one fixed, tiny, fully worked example: a $K=2$ expert model at a
single fixed input $\bm x$ (we drop the $\bm x$ argument since it never changes), with
\[
w_1 = 0.6,\quad \mu_1 = 2,\quad \sigma_1 = 0.5, \qquad\qquad
w_2 = 0.4,\quad \mu_2 = 6,\quad \sigma_2 = 1.0.
\]
By Eq.~\eqref{eq:yhatdef}, $\hat y = 0.6\cdot 2 + 0.4\cdot 6 = 1.2 + 2.4 = 3.6$. By
Eq.~\eqref{eq:lotv},
\[
\sigma_a^2 = 0.6\cdot 0.5^2 + 0.4\cdot 1.0^2 = 0.15 + 0.40 = 0.55,\qquad
\sigma_e^2 = 0.6\,(2-3.6)^2 + 0.4\,(6-3.6)^2 = 0.6\cdot2.56 + 0.4\cdot5.76 = 1.536+2.304 = 3.84,
\]
so $\sigma_{\mathrm{tot}}^2 = 0.55+3.84 = 4.39$, i.e.\ $\sigma_{\mathrm{tot}} \approx 2.095$.
Notice that the epistemic term $(3.84)$ dwarfs the aleatoric term $(0.55)$ here: almost all of
this point's total uncertainty comes from the two experts disagreeing about \emph{where} the
outcome is (2 vs.\ 6), not from either expert individually being unsure. This is exactly the
regime -- large expert disagreement -- where the ideas in this document matter most, and it is
why Section~\ref{sec:seta} eventually introduces a knob that acts on $\sigma_e^2$ specifically.
\end{example}

We now leave the model fixed (Example~\ref{ex:running}'s two numbers never change) and build up,
step by step, everything needed to turn $(\mu_k,\sigma_k,w_k)$ into a valid, tight prediction set.

% ============================================================================
\section{Conformal prediction, from scratch}
\label{sec:cp}

This section builds \emph{split conformal prediction} \cite{vovk2005algorithmic,lei2018distribution},
the machinery that gives \emph{any} prediction set the coverage guarantee of
Eq.~\eqref{eq:coverage-def}, exactly, for any model whatsoever -- including a badly wrong one. It
is the part of the story that makes coverage a solved problem, so that everything after this
section (Sections~\ref{sec:hpd}--\ref{sec:seta}) is free to focus entirely on width.

\subsection{Setup and the one assumption: exchangeability}

We have a trained model (Section~\ref{sec:moe}) that we are not going to retrain or modify. We
additionally have a \emph{calibration set} $\mathcal D_{\mathrm{cal}} = \{(\bm x_i, y_i)\}_{i=1}^n$
-- labeled data the model was not trained on -- and we will eventually see a fresh test point
$(\bm x_{\mathrm{test}}, y_{\mathrm{test}})$ whose label we do not get to see, only predict a set
for.

The one assumption every method in this document relies on is \emph{exchangeability}: the $n+1$
points $(\bm x_1,y_1),\dots,(\bm x_n,y_n),(\bm x_{\mathrm{test}},y_{\mathrm{test}})$, taken
together, have a joint distribution that does not change if we relabel/permute which one is
``the test point.'' (Independent, identically-distributed data is one common way to get
exchangeability, but exchangeability is a strictly weaker requirement.) Nothing here assumes the
trained model is correct, that the errors are Gaussian, or that $y$ depends on $\bm x$ in any
particular way -- only this one symmetry condition on how the data was generated.

\subsection{Nonconformity scores}

Fix any rule for turning a labeled point $(\bm x, y)$ into a single real number $S(\bm x, y)$,
called a \emph{nonconformity score}: intuitively, ``how surprising is it that this $\bm x$ was
paired with this $y$, according to my model?'' A large score means ``very surprising, doesn't fit
my model's story for $\bm x$ well''; a small score means ``exactly what I'd expect.'' The only
formal requirement on $S$ for what follows is that it is a fixed, deterministic function -- it
does not get to look at, or be re-optimized using, the label of the point currently being scored,
beyond the specific mechanics described in Section~\ref{sec:validity} later. Sections
\ref{sec:standardcp}--\ref{sec:seta} are all, structurally, ``choose a different $S$'' with
everything else in this section held fixed.

\subsection{The finite-sample quantile correction}
\label{sec:quantile}

Compute the score $S_i = S(\bm x_i, y_i)$ for every calibration point, giving $n$ numbers
$S_1,\dots,S_n$. Sort them ascending: $S_{(1)} \le S_{(2)} \le \dots \le S_{(n)}$. The threshold we
use is \emph{not} simply the empirical $(1-\alpha)$ quantile (e.g.\ the $90$th percentile for
$\alpha=0.1$); it is a slightly more conservative, finite-sample-corrected version:
\begin{equation}
\hat q = S_{(k)}, \qquad k = \Big\lceil (n+1)(1-\alpha) \Big\rceil,
\label{eq:qlevel-order}
\end{equation}
with the convention that if $k>n$ (which happens for small $n$ or $\alpha$), $\hat q = +\infty$
(the resulting prediction set is the whole real line -- coverage can trivially be met, at the cost
of an unusable, infinitely wide set; this only occurs for tiny calibration sets). Equivalently, in
terms of a quantile \emph{level} rather than an order-statistic index,
\begin{equation}
q_{\mathrm{level}} = \min\!\left(1,\ \frac{\lceil (n+1)(1-\alpha)\rceil}{n}\right), \qquad
\hat q = \operatorname{Quantile}_{q_{\mathrm{level}}}(\{S_i\}),
\label{eq:qlevel}
\end{equation}
using the ``upper'' (round up to the next actually-observed data point) quantile convention. Every
calibrator described in this document uses exactly this construction; only the score $S$ changes.

\begin{example}
With $n=19$ calibration points and $\alpha=0.1$: $k = \lceil 20\times0.9\rceil = \lceil 18\rceil =
18$, so $\hat q$ is the \emph{18th smallest} of the 19 scores -- not the 90th percentile in the
everyday sense (which would sit between the 17th and 18th of 19 sorted points), but pushed one
notch higher, to the largest value below the very worst score. This one-notch conservatism is
exactly what Theorem~\ref{thm:coverage} below shows is necessary and sufficient for exact
coverage.
\end{example}

\subsection{Why this gives exact coverage}
\label{sec:validity}

\begin{theorem}[Split-conformal marginal coverage]
\label{thm:coverage}
Suppose $S$ is fixed independently of the calibration and test labels (in particular, not tuned
using $y_1,\dots,y_n$ or $y_{\mathrm{test}}$), and the $n+1$ points are exchangeable. Then, with
$\hat q$ as in Eq.~\eqref{eq:qlevel-order} and $\hat C(\bm x_{\mathrm{test}}) = \{y :
S(\bm x_{\mathrm{test}}, y) \le \hat q\}$,
\begin{equation}
P\big(y_{\mathrm{test}} \in \hat C(\bm x_{\mathrm{test}})\big) \ge 1-\alpha.
\end{equation}
\end{theorem}
\begin{proof}
Let $S_{\mathrm{test}} = S(\bm x_{\mathrm{test}}, y_{\mathrm{test}})$ and consider the $n+1$
numbers $S_1,\dots,S_n,S_{\mathrm{test}}$. Because $S$ does not depend on which point is labeled
``test'' and the underlying $n+1$ points are exchangeable, these $n+1$ scores are themselves
exchangeable, so (ties broken arbitrarily) the \emph{rank} of $S_{\mathrm{test}}$ among all $n+1$
scores is uniformly distributed on $\{1,\dots,n+1\}$: every position is equally likely a priori,
by symmetry. The event $y_{\mathrm{test}} \in \hat C(\bm x_{\mathrm{test}})$ is exactly the event
$S_{\mathrm{test}} \le \hat q = S_{(k)}$, i.e.\ that $S_{\mathrm{test}}$'s rank among the $n+1$
combined scores is at most $k$. By uniformity of the rank,
$P(\text{rank}(S_{\mathrm{test}}) \le k) = k/(n+1) \ge (1-\alpha)$, using $k = \lceil
(n+1)(1-\alpha)\rceil \ge (n+1)(1-\alpha)$.
\end{proof}

This is the single most important fact in this entire document: \emph{coverage is free}, for any
choice of nonconformity score $S$ (subject only to the ``fixed in advance of the labels''
condition), regardless of whether the underlying model is any good. Every method from here on
picks a specific, cleverer $S$, purely to make the resulting set \emph{narrower}, never to make it
\emph{more valid} -- validity is already guaranteed by Theorem~\ref{thm:coverage} the moment $S$
is fixed.

\subsection{Two classical calibrators, worked through in full}
\label{sec:standardcp}

To make Section~\ref{sec:quantile}--\ref{sec:validity} completely concrete before anything more
sophisticated is introduced, here are the two simplest possible choices of $S$.

\paragraph{Standard CP.} Ignore every uncertainty estimate the model reports and score the raw
absolute residual, $S_i = |y_i - \hat y(\bm x_i)|$. The threshold $\hat q$ from
Eq.~\eqref{eq:qlevel-order} is then a single number, and the resulting prediction interval is
$\hat y(\bm x) \pm \hat q$ -- the same width for every test point, regardless of how confident or
unconfident the model is at $\bm x$. This is valid (Theorem~\ref{thm:coverage} always applies) but
wasteful: it must be wide enough to cover the model's worst, least-confident predictions, and pays
that same width even on inputs where the model is highly confident.

\paragraph{CP-VS (variance-scaled CP).} Normalize the residual by the model's own reported total
scale before scoring, $S_i = |y_i - \hat y(\bm x_i)| / \sigma_{\mathrm{tot}}(\bm x_i)$. The
resulting interval is $\hat y(\bm x) \pm \hat q\, \sigma_{\mathrm{tot}}(\bm x)$: heteroscedastic,
wider exactly where the model itself reports more uncertainty. This is a strict improvement over
Standard CP whenever $\sigma_{\mathrm{tot}}(\bm x)$ is genuinely informative about $|y-\hat
y(\bm x)|$, and it is still exactly valid by Theorem~\ref{thm:coverage} -- normalizing by
$\sigma_{\mathrm{tot}}(\bm x_i)$ is just another fixed-in-advance scoring rule.

Both of these commit, by construction, to a \emph{symmetric interval} centered at $\hat y(\bm x)$.
A third classical alternative, conformalized quantile regression (CQR) \cite{romano2019conformalized},
instead conformalizes an auxiliary quantile-regression head and shifts its predicted quantiles by a
single additive correction -- still not shaped by the model's own predictive density. Section~\ref{sec:hpd}
shows why, once that density is genuinely multimodal (as in Example~\ref{ex:running}), ignoring its
shape is itself a source of wasted width -- independent of how well-tuned any single scale
$\sigma_{\mathrm{tot}}(\bm x)$ is.

% ============================================================================
\section{Highest-density regions: the right shape for the set}
\label{sec:hpd}

\subsection{Why a symmetric interval can be the wrong shape}

Look again at Example~\ref{ex:running}. If we imagine, just for a moment, that the model's
mixture density $f_{\mathrm{mix}}(\cdot\mid \bm x)$ from Eq.~\eqref{eq:mixdef} \emph{were} exactly
the true conditional density of $y$ (we will make this exact and remove the ``imagine'' in
Section~\ref{sec:moghpd}), plug in the two components:
\[
f_{\mathrm{mix}}(y) = 0.6\,\mathcal N(y; 2, 0.5^2) + 0.4\,\mathcal N(y; 6, 1.0^2).
\]
Evaluate this density at three points: the two expert means, and the aggregate prediction $\hat y
= 3.6$ that a symmetric interval would be centered on.
\begin{align*}
f_{\mathrm{mix}}(2) &= 0.6\cdot \tfrac{1}{0.5\sqrt{2\pi}}\cdot 1 \ +\ 0.4\cdot(\text{negligible})
\ \approx\ 0.4788,\\
f_{\mathrm{mix}}(6) &= 0.4\cdot \tfrac{1}{1.0\sqrt{2\pi}}\cdot 1 \ +\ (\text{negligible})
\ \approx\ 0.1596,\\
f_{\mathrm{mix}}(3.6) &\approx\ 0.6\cdot0.004768 \ +\ 0.4\cdot0.02239 \ \approx\ 0.0118
\end{align*}
(numerically: at $y=3.6$, expert 1 is $3.2$ of its own standard deviations away and expert 2 is
$2.4$ of its own standard deviations away, so both components have already decayed a great deal).
The density at $\hat y=3.6$ -- the exact center any symmetric interval would use -- is more than
$40\times$ smaller than the density at either mode. A symmetric interval $\hat y \pm s$ must, to
reach any reasonable coverage, either stretch $s$ far enough to swallow the near-empty valley
between the two modes (wasting width on values of $y$ that are individually very unlikely), or
fail to reach one of the two modes entirely. Neither option is necessary: the region that actually
achieves a given coverage with the least total width is not an interval at all here, but
\emph{two} disjoint intervals, one bracketing each mode, with nothing in between. This is the
motivating picture for everything in this section.

\subsection{Definition}

\begin{definition}[Highest-density region]
\label{def:hpd}
For a probability density $f$ on $\mathbb R$ and $\alpha\in(0,1)$, the $100(1-\alpha)\%$
\emph{highest-density region} of $f$ is
\begin{equation}
R_\alpha(f) = \{y \in \mathbb R : f(y) \ge \tau_\alpha\},
\end{equation}
where $\tau_\alpha$ is the largest constant such that $\int_{R_\alpha(f)} f \ge 1-\alpha$. In
words: keep exactly the values of $y$ where the density is \emph{tallest}, stopping once you have
accumulated $1-\alpha$ of the total probability mass. Depending on $f$'s shape, $R_\alpha(f)$ can
be a single interval (unimodal $f$) or a union of several disjoint intervals (multimodal $f$).
\end{definition}

This is a completely classical, pre-conformal-prediction idea from Bayesian and frequentist
density estimation, originally studied as the natural generalization of a symmetric credible
interval to arbitrary densities \cite{boxtiao1973bayesian,hyndman1996computing}. The next theorem
is a fact purely about a fixed density $f$ -- it has nothing to do with calibration sets, models,
or conformal prediction yet.

\subsection{Why the level set is width-optimal}

\begin{theorem}[Level-set width optimality]
\label{thm:hpd}
Fix a density $f$ and $\alpha\in(0,1)$, and let $R_\tau=\{y:f(y)\ge\tau\}$ with $\tau$ chosen so
$\int_{R_\tau}f=1-\alpha$ (i.e.\ $R_\tau=R_\alpha(f)$). Then for \emph{any} measurable set $C$
with $\int_C f \ge 1-\alpha$, $|R_\tau| \le |C|$, where $|\cdot|$ denotes total length (Lebesgue
measure).
\end{theorem}

\emph{Intuition before the proof.} Think of $f$ as describing how ``valuable'' (in probability
mass per unit length) each point $y$ is. $R_\tau$ greedily buys the most valuable points first,
stopping once it has bought $1-\alpha$ of the total value. Any other set $C$ with the same total
value must, compared to $R_\tau$, be missing some of $R_\tau$'s points (each worth at least
$\tau$ per unit length) and including some points $R_\tau$ left out (each worth less than $\tau$
per unit length, or $C$ could not need more length than $R_\tau$ to hit the same value). Trading
away high-value length for low-value length to reach the same total value always costs you extra
length. That is the entire proof, made precise:

\begin{proof}
Let $D_1 = R_\tau \setminus C$ and $D_2 = C\setminus R_\tau$. Since $D_1 \subset R_\tau$, $f\ge
\tau$ on $D_1$; since $D_2 \subset R_\tau^c$, $f < \tau$ on $D_2$. The region $R_\tau \cap C$ is
shared by both sets, so $\int_C f \ge \int_{R_\tau} f = 1-\alpha \ge \int_C f$ forces
$\int_{D_2} f \ge \int_{D_1} f$ (the extra mass $C$ has over $R_\tau\cap C$, from $D_2$, must be at
least as large as the mass $R_\tau$ has over $R_\tau\cap C$, from $D_1$, since both totals equal
$1-\alpha$ once you add back $\int_{R_\tau\cap C} f$ to each). Hence
$\tau|D_2| \ge \int_{D_2}f \ge \int_{D_1}f \ge \tau|D_1|$, so $|D_2|\ge|D_1|$, and
$|R_\tau|-|C| = |D_1|-|D_2| \le 0$.
\end{proof}

\begin{corollary}[A single Gaussian recovers the symmetric interval]
\label{cor:gauss-hpd}
If $f=\mathcal N(\mu,\sigma^2)$ (a single, unimodal, symmetric Gaussian), $R_\alpha(f)$ is exactly
the familiar symmetric interval $[\mu - \sigma z_{1-\alpha/2},\ \mu+\sigma z_{1-\alpha/2}]$, where
$z_{1-\alpha/2}$ is the standard-normal $(1-\alpha/2)$ quantile.
\end{corollary}
\begin{proof}
$f$ is symmetric and strictly decreasing in $|y-\mu|$, so every super-level set $\{f\ge\tau\}$ is a
single interval centered at $\mu$; choosing its half-width so it carries mass $1-\alpha$ under a
standard Gaussian gives $\sigma z_{1-\alpha/2}$ exactly.
\end{proof}

Corollary~\ref{cor:gauss-hpd} is the reassurance that highest-density regions are not some exotic
replacement for the familiar symmetric interval -- they \emph{are} the familiar symmetric interval,
in the one case (a single unimodal, symmetric density) where the symmetric interval is already
correct. Section~\ref{sec:standardcp}'s Standard CP and CP-VS are implicitly always using the
\emph{equal-tailed} region $[F^{-1}(\alpha/2), F^{-1}(1-\alpha/2)]$ instead (excluding equal
probability mass from each tail); by Theorem~\ref{thm:hpd} this can only be as narrow as
$R_\alpha(f)$, with equality exactly in the unimodal-symmetric case of
Corollary~\ref{cor:gauss-hpd}. For a genuinely multimodal $f$, like Example~\ref{ex:running}'s
mixture, the two constructions diverge, and $R_\alpha(f)$ is strictly narrower.

\subsection{MoG-HPD: conformalizing the level set}
\label{sec:moghpd}

Theorem~\ref{thm:hpd} is a fact about a \emph{known} density $f$. In practice we do not know the
true conditional density of $y\mid\bm x$; all we have is the model's fitted mixture
$f_{\mathrm{mix}}(\cdot\mid\bm x)$ of Eq.~\eqref{eq:mixdef}, which is at best an approximation to
the truth. \textbf{MoG-HPD} \cite{izbicki2020cd,sadinle2019least} is the calibrator that plugs
$f=f_{\mathrm{mix}}(\cdot\mid\bm x)$ into
Definition~\ref{def:hpd} and then uses Section~\ref{sec:cp}'s machinery to pick the threshold
$\tau$ in a way that is guaranteed valid \emph{regardless of whether $f_{\mathrm{mix}}$ is
correct}. Concretely: the nonconformity score is the negative log-density,
\begin{equation}
S_i = -\log f_{\mathrm{mix}}(y_i \mid \bm x_i),
\end{equation}
conformalized exactly as in Section~\ref{sec:quantile} to get a threshold $\hat t = \hat q$ (using
the notation $\hat t$ from here on, since it lives on the $-\log$ scale), and the prediction set
returned for a new $\bm x$ is
\begin{equation}
\hat C(\bm x) = \big\{y : -\log f_{\mathrm{mix}}(y\mid\bm x) \le \hat t\big\}
= \big\{y : f_{\mathrm{mix}}(y\mid\bm x) \ge e^{-\hat t}\big\},
\label{eq:moghpd-region}
\end{equation}
i.e.\ exactly the empirical highest-density region of $f_{\mathrm{mix}}(\cdot\mid\bm x)$ at level
$\hat\tau = e^{-\hat t}$. Two separate facts hold about Eq.~\eqref{eq:moghpd-region}, and it is
important to keep them apart:
\begin{itemize}[leftmargin=1.5em]
\item \textbf{Validity} follows immediately from Theorem~\ref{thm:coverage}: $S_i = -\log
f_{\mathrm{mix}}(y_i\mid\bm x_i)$ is a fixed function of $(\bm x_i,y_i)$ (it does not peek at
$y_{\mathrm{test}}$), so $P(y_{\mathrm{test}}\in\hat C(\bm x_{\mathrm{test}}))\ge 1-\alpha$
\emph{regardless of whether $f_{\mathrm{mix}}$ is a good model of the truth}.
\item \textbf{Width-optimality} only holds \emph{relative to $f_{\mathrm{mix}}$ itself}
(Theorem~\ref{thm:hpd}): if the fitted mixture is close to the true conditional density, the
resulting set is close to the truly optimal set; if $f_{\mathrm{mix}}$ is systematically wrong in
some way, MoG-HPD faithfully returns the level set of the \emph{wrong} density -- valid, but not
as narrow as it could be. This is exactly the gap that Sections~\ref{sec:stahpd} and
\ref{sec:seta} exist to close.
\end{itemize}

\subsubsection{Finding the region numerically}
\label{sec:numerics}

Equation~\eqref{eq:moghpd-region} defines the region implicitly (as wherever a curve clears a
threshold), not as a closed-form interval, so it has to be located numerically. This repository's
implementation (\texttt{MoG\_HPD\_Calibrator.\_segments\_for\_point}) does this in three steps,
which every later method (STA-HPD, SETA-HPD) reuses unchanged:
\begin{enumerate}[leftmargin=1.5em]
\item \textbf{Window.} Build a search window wide enough to contain all the mass, starting from
$[\min_k \mu_k - 6\sigma_k,\ \max_k \mu_k + 6\sigma_k]$ and doubling outward (up to a fixed number
of times) if the density at the window's edges has not yet decayed to a negligible fraction of the
threshold.
\item \textbf{Grid scan.} Evaluate $f_{\mathrm{mix}}(y) - \tau$ on a fine grid inside the window
(finer near each component's own mean, to avoid missing a narrow peak) and record every pair of
adjacent grid points where the sign flips -- each such pair brackets one boundary of the region.
\item \textbf{Root refinement.} Within each bracketing pair, run Brent's method (a fast, reliable
1-D root finder that only needs the function values, not a derivative) to pin the exact boundary
to numerical precision. The boundaries, sorted and paired up, are the disjoint $[{\rm lower},{\rm
upper}]$ segments of $\hat C(\bm x)$; segments closer together than a small numerical tolerance
are merged.
\end{enumerate}
The returned prediction set for a test point is therefore not always a single interval: it is a
\emph{ragged} list of $[\text{lower},\text{upper}]$ segments, possibly empty (rare, only under
numerical edge cases), possibly a single interval (the common case, and the only possible case
when $K=1$), and possibly several disjoint intervals (when the mixture is genuinely multimodal and
the threshold falls in a valley, exactly as in Example~\ref{ex:running}).

\begin{example}[Example~\ref{ex:running}, continued]
Suppose, hypothetically, that a calibration procedure had produced the threshold $\hat\tau =
0.05$ for this model (this is smaller than $f_{\mathrm{mix}}(3.6)\approx0.0118$ from earlier --
wait, we need $\hat\tau$ \emph{above} the valley for it to open a gap; take instead $\hat\tau =
0.05$, which sits between the valley height $0.0118$ and the shorter peak height $f_{\mathrm{mix}}(6)
\approx 0.1596$). Since $f_{\mathrm{mix}}(3.6) = 0.0118 < \hat\tau = 0.05$, the region $\{y :
f_{\mathrm{mix}}(y)\ge\hat\tau\}$ does \emph{not} include $y=3.6$: it is two disjoint intervals,
one bracketing $y=2$ (where the tall, narrow first component clears $0.05$) and one bracketing
$y=6$ (where the shorter, wider second component clears $0.05$), with a gap in between covering
the near-empty valley. A symmetric interval centered at $\hat y=3.6$ could never produce this
shape, no matter how its half-width $s$ is tuned -- it would have to either include the whole
valley or miss one of the two modes.
\end{example}

% ============================================================================
\section{STA-HPD: correcting how the density is scaled and thresholded}
\label{sec:stahpd}

\subsection{What is still wrong with plain HPD calibration}

MoG-HPD (Section~\ref{sec:moghpd}) is valid by construction and shape-optimal \emph{relative to
$f_{\mathrm{mix}}$}. But $f_{\mathrm{mix}}$ is a trained neural network's output, not the truth,
and two specific, systematic mismatches recur across real datasets:
\begin{enumerate}[leftmargin=1.5em]
\item \textbf{Global over/under-dispersion.} The per-expert standard deviations $\sigma_k(\bm x)$
might be systematically too large or too small relative to how big the actual residuals turn out
to be, uniformly across the dataset -- e.g., a model trained with a particular loss weighting might
consistently over-report or under-report its own spread.
\item \textbf{Wrong sensitivity to the model's own scale.} Even when $\sigma_k(\bm x)$ is centered
correctly on average, it may carry more or less \emph{information} about the actual residual size
on some datasets than others (a large reported $\sigma$ does not always mean a large error is
likely, if $\sigma$ itself is a noisy quantity) -- so trusting it with a fixed, dataset-independent
curvature (as MoG-HPD implicitly does) is not always right.
\end{enumerate}
\textbf{STA-HPD} (Shape--Threshold Adaptive HPD) opens exactly two tunable knobs to correct these,
each a no-op at its default value, so that MoG-HPD is always reachable as a special case.

\subsection{The shape exponent $c$}
\label{sec:c}

Freeze, once, on the full calibration set, the geometric mean of every per-expert reported
standard deviation,
\begin{equation}
G = \exp\!\Big(\operatorname*{mean}_{k,i}\ \log\sigma_k(\bm x_i)\Big),
\label{eq:Gdef}
\end{equation}
i.e.\ the ``typical'' scale this model reports across the whole calibration set. For a
\emph{shape} exponent $c\in\mathbb R_{\ge0}$, define the reshaped per-expert scale
\begin{equation}
\tilde\sigma_k(\bm x; c) = G\,\big(\sigma_k(\bm x)/G\big)^c.
\label{eq:sigmashape}
\end{equation}

\begin{lemma}[What $c$ actually does]
\label{lem:cdoes}
$\log\tilde\sigma_k(\bm x;c) = \log G + c\big(\log\sigma_k(\bm x)-\log G\big)$, i.e.\ reshaping is
a straight line in log-scale space, anchored at $G$, with slope $c$. Consequently
$\tilde\sigma_k(\bm x;1)=\sigma_k(\bm x)$ (the raw mixture, unchanged); $\tilde\sigma_k(\bm x;0)=G$
for every $k$ and $\bm x$ (every expert's scale collapses to the same constant $G$, i.e.\
heteroscedasticity is switched off entirely); and for $c\in(0,1)$, every $\sigma_k(\bm x)$ is
pulled part-way toward $G$ (large $\sigma_k$ shrunk down, small $\sigma_k$ pulled up), while $c>1$
does the opposite, stretching the spread of scales apart.
\end{lemma}
\begin{proof}
Take logs of Eq.~\eqref{eq:sigmashape}. The three boundary facts follow by substituting $c=1$,
$c=0$, and reading off the slope $c$ of the resulting affine function.
\end{proof}

\begin{example}[Example~\ref{ex:running}, continued]
With only these two numbers on the ``calibration set,'' $G = \exp\!\big(\tfrac12(\log 0.5 +
\log 1.0)\big) = \exp(-0.3466) \approx 0.7071$ (the geometric mean of $0.5$ and $1.0$, which is
$\sqrt{0.5\times1.0}=\sqrt{0.5}\approx0.7071$). At $c=0.5$:
\[
\tilde\sigma_1(0.5) = 0.7071\cdot(0.5/0.7071)^{0.5} = 0.7071\times0.8409 \approx 0.5946, \qquad
\tilde\sigma_2(0.5) = 0.7071\cdot(1.0/0.7071)^{0.5} = 0.7071\times1.1892 \approx 0.8409.
\]
The narrower expert's scale moved \emph{up} (from $0.5$ toward $0.7071$) and the wider expert's
scale moved \emph{down} (from $1.0$ toward $0.7071$): $c<1$ genuinely compresses the two experts'
reported scales toward their common geometric-mean anchor, exactly as Lemma~\ref{lem:cdoes}
predicts.
\end{example}

\subsection{The threshold-field exponent $\theta$}
\label{sec:theta}

Define the reshaped mixture density and its total predictive scale (the aleatoric term now uses
the reshaped $\tilde\sigma_k$; the epistemic term is untouched -- $c$ is designed to act on the
within-expert spread specifically, per failure mode 1 above),
\begin{align}
f_c(y\mid\bm x) &= \sum_{k=1}^K w_k(\bm x)\,\mathcal N\big(y;\ \mu_k(\bm x),\
\tilde\sigma_k^2(\bm x;c)\big), \label{eq:fc}\\
\sigma_{\mathrm{tot}}^2(\bm x;c) &= \sum_{k=1}^K w_k(\bm x)\,\tilde\sigma_k^2(\bm x;c) \;+\;
\sum_{k=1}^K w_k(\bm x)\,\big(\mu_k(\bm x)-\hat y(\bm x)\big)^2. \label{eq:sigmatotc}
\end{align}
A \emph{second}, separate exponent $\theta \in \mathbb R$ then tilts the conformal score itself by
the point's own total scale, correcting failure mode 2 above:
\begin{equation}
S_i(c,\theta) = -\log f_c(y_i\mid\bm x_i) \;-\; \theta\,\log\sigma_{\mathrm{tot}}(\bm x_i;c),
\label{eq:stascore}
\end{equation}
with threshold $\hat t(c,\theta)$ from Section~\ref{sec:quantile} applied to these scores, and
induced region
\begin{equation}
\hat C(\bm x) = \Big\{y :\ \log f_c(y\mid\bm x) \;\ge\; -\hat t(c,\theta) \;-\;
\theta\log\sigma_{\mathrm{tot}}(\bm x;c) \Big\}.
\label{eq:staregion}
\end{equation}
Read Eq.~\eqref{eq:staregion} carefully: it is \emph{still} a level set of the single density
$f_c$ (the same object at every $\bm x$, once $c$ is fixed), but the \emph{threshold it is compared
against} is no longer the single global $e^{-\hat t}$ that MoG-HPD used -- it is scaled up or down,
per test point, by that point's own $\sigma_{\mathrm{tot}}(\bm x;c)^{-\theta}$. $\theta>0$ demands
\emph{more} density (a stricter threshold, narrower set, all else equal) exactly where the model
reports more total scale; $\theta<0$ does the opposite; $\theta=0$ recovers the single global
threshold.

\subsection{Boundary recovery: STA-HPD contains everything introduced so far}

\begin{proposition}[Boundary recovery]
\label{prop:starecovery}
$(c,\theta)=(1,0)$ recovers MoG-HPD exactly (Eqs.~\ref{eq:stascore}--\ref{eq:staregion} reduce to
Eq.~\ref{eq:moghpd-region} directly). When $K=1$ (a single Gaussian, no mixture at all),
$(c,\theta)=(1,1)$ recovers CP-VS (Section~\ref{sec:standardcp}) on the total predictive standard
deviation exactly. $(c,\theta)=(0,0)$ recovers Standard CP: every $\tilde\sigma_k(\bm x;0)=G$
collapses to the same constant (Lemma~\ref{lem:cdoes}), so $f_0(\cdot\mid\bm x)$ no longer depends
on $\bm x$ through the experts' scales at all, and $\theta=0$ leaves a single global threshold, so
the whole construction reduces to a single fixed-width region.
\end{proposition}
\begin{proof}
Direct substitution into Eqs.~\eqref{eq:sigmashape}--\eqref{eq:staregion}, using
Lemma~\ref{lem:cdoes}'s two boundary identities $\tilde\sigma_k(\bm x;1)=\sigma_k(\bm x)$ and
$\tilde\sigma_k(\bm x;0)=G$.
\end{proof}

\begin{corollary}[Closed form when $K=1$]
\label{cor:k1closedform}
When $K=1$, write $\tilde\sigma(\bm x;c)$ for the sole reshaped scale and $\hat\tau=e^{-\hat
t(c,\theta)}$. The region $\hat C(\bm x)$ of Eq.~\eqref{eq:staregion} is the symmetric interval
$\hat y(\bm x) \pm s(\bm x)$ with
\begin{equation}
s(\bm x) = \tilde\sigma(\bm x;c)\ \sqrt{2\log\!\left(\frac{1}{\hat\tau\sqrt{2\pi}\
\tilde\sigma(\bm x;c)^{1-\theta}}\right)}.
\label{eq:k1closedform}
\end{equation}
\end{corollary}
\begin{proof}
For $K=1$, $f_c(y\mid\bm x)$ is a single Gaussian with scale $\tilde\sigma(\bm x;c)$, and
$\sigma_{\mathrm{tot}}(\bm x;c)=\tilde\sigma(\bm x;c)$ exactly (the epistemic term is identically
zero with only one expert). Substituting the Gaussian density into $\log f_c(y\mid\bm x)\ge -\hat
t-\theta\log\tilde\sigma(\bm x;c)$ and solving for $|y-\mu(\bm x)|$ gives
Eq.~\eqref{eq:k1closedform} after collecting terms and taking a square root.
\end{proof}

This closed form is a useful sanity check: plugging $(c,\theta)=(1,0)$ into
Eq.~\eqref{eq:k1closedform} gives back MoG-HPD's single-Gaussian half-width exactly; plugging in
$(1,1)$ makes the $\tilde\sigma(\bm x;c)^{1-\theta}$ factor vanish, giving $s(\bm x) =
\hat q\,\sigma(\bm x)$ for a constant $\hat q$ -- CP-VS's form exactly.

\subsection{Why this can actually help: a population-level argument}
\label{sec:optimality}

Section~\ref{sec:hpd} established that a level set is width-optimal \emph{for the density it is a
level set of}. It says nothing about which density -- or which threshold field -- is the right
one to use in the first place. This subsection derives what the truly optimal choice would have to
satisfy, so that the gap between that ideal and what plain MoG-HPD does is visible, and so that
$(c,\theta)$'s role in closing that gap is not just a heuristic.

\begin{theorem}[First-order optimality of the level-set family]
\label{thm:optimality}
Let the true conditional density of $Y\mid\bm x$ be $p(\cdot\mid\bm x)$ (generally unknown, and
generally not equal to any $f_c$). Consider choosing a per-point log-threshold function $g(\bm x)$,
defining $C_g(\bm x) = \{y : \log f(y\mid\bm x) \ge -g(\bm x)\}$ for whatever fixed density $f$ is
in force, to solve
\[
\min_g\ \mathbb E_{\bm x}\big[|C_g(\bm x)|\big] \quad\text{subject to}\quad
\mathbb E_{\bm x}\big[P(Y\in C_g(\bm x)\mid \bm x)\big] \ge 1-\alpha.
\]
At a stationary point of the associated Lagrangian, writing $b(\bm x)$ for a boundary point of
$C_g(\bm x)$,
\begin{equation}
p\big(b(\bm x)\mid\bm x\big) = \text{a constant, the same for every } \bm x.
\label{eq:optcond}
\end{equation}
\end{theorem}
\emph{Intuition.} If the true density at the boundary were higher at some $\bm x$ than at another,
you could shave a sliver of width off the low-boundary-density point (losing little true coverage
there, since the true density is low there) and spend it at the high-boundary-density point
(buying more true coverage there per unit of width spent) -- a strict improvement, until the
boundary densities equalize. The full derivation:
\begin{proof}
\emph{Step 1 (boundary perturbation).} Perturb $g(\bm x)\mapsto g(\bm x)+\delta$ at one $\bm x$
only. Each boundary point $b(\bm x)$ solves $\log f(b(\bm x)\mid\bm x)=-g(\bm x)$; linearizing,
$\Delta b = -\delta\, f(b\mid\bm x)/f'(b\mid\bm x)$ to first order.

\emph{Step 2 (width response).} The set's width changes, to first order, by
$\delta\, f(b\mid\bm x)/|f'(b\mid\bm x)|$ (summed over boundary points, with sign chosen so
increasing $\delta$ enlarges the region).

\emph{Step 3 (true-coverage response).} The \emph{true} covered probability
$\int_{C_g(\bm x)} p(y\mid\bm x)\,dy$ changes, by the same boundary displacement, by $\delta\,
p(b\mid\bm x)\, f(b\mid\bm x)/|f'(b\mid\bm x)|$ -- identical to Step 2's rate except the model
density $f(b\mid\bm x)$ in the numerator is replaced by the \emph{true} density $p(b\mid\bm x)$.

\emph{Step 4 (stationarity).} The Lagrangian is $\mathcal L[g] = \mathbb E_{\bm x}[|C_g(\bm x)|] -
\lambda\big(\mathbb E_{\bm x}[P(Y\in C_g(\bm x)\mid\bm x)] - (1-\alpha)\big)$. Stationarity requires
Step 2's rate to equal $\lambda\times$ Step 3's rate at every $\bm x$. The common factor
$f(b\mid\bm x)/|f'(b\mid\bm x)|$ cancels from both sides, leaving $1 = \lambda\, p(b(\bm x)\mid\bm
x)$ at every $\bm x$, for one shared multiplier $\lambda$ -- i.e.\ $p(b(\bm x)\mid\bm x) =
1/\lambda$, the same value everywhere.
\end{proof}

Plain MoG-HPD instead enforces $f_{\mathrm{mix}}(b(\bm x)\mid\bm x) = \hat\tau$, a constant
condition on the \emph{model's} density -- which coincides with the true optimum
(Eq.~\ref{eq:optcond}) only where $f_{\mathrm{mix}}$ agrees with the truth exactly at the
boundary. The shape exponent $c$ corrects the part of that mismatch driven by systematic
mis-scaling of $\sigma_k(\bm x)$; the threshold field $\theta$ corrects the remaining pointwise
level mismatch. In particular, when $\sigma_k(\bm x)$ only partially tracks the true residual size,
dispersion in $\sigma_k$ across the dataset inflates typical values of $\sigma_k^{c}$-type terms
faster than it deflates the matching conformal quantile, so the width-minimizing $c^\star$ is
generally interior and can fall strictly below $1$ -- a data-driven shrinkage of heteroscedasticity
that plain MoG-HPD (which fixes $c\equiv1$) cannot express.

\subsection{Fitting $(c,\theta)$ without breaking validity}
\label{sec:stafit}

There is a subtlety here that is easy to get wrong. Theorem~\ref{thm:coverage} requires $S$ to be
fixed \emph{independently of the labels it is being conformalized against}. If $(c,\theta)$ were
chosen by minimizing width on the very same calibration points that then supply $\hat t(c,\theta)$,
the resulting score would have secretly used the labels $y_1,\dots,y_n$ twice -- once to pick
$(c,\theta)$, once to pick $\hat t$ -- breaking exchangeability and, with it,
Theorem~\ref{thm:coverage}'s guarantee.

The fix is to split $\mathcal D_{\mathrm{cal}}$ into two disjoint pieces, $\mathcal
D_{\mathrm{tune}}$ (used only to \emph{choose} $(c,\theta)$) and $\mathcal D_{\mathrm{calib}}$
(used only to compute the final $\hat t$, at the now-fixed $(c,\theta)$). Because the split is
made before either fold's labels are used for anything, and $\mathcal D_{\mathrm{calib}}$'s
labels are never touched during the $(c,\theta)$ search, the final score is a fixed function of
$(\bm x,y)$ by the time it is conformalized, and Theorem~\ref{thm:coverage} applies exactly, with
no assumption about model correctness. This is stated precisely below.

\begin{proposition}[Marginal validity]
\label{prop:stavalidity}
For any $(c,\theta)$ selected using only $\mathcal D_{\mathrm{tune}}$, disjoint from the $\mathcal
D_{\mathrm{calib}}$ fold that supplies $\hat t$, Eq.~\eqref{eq:stascore} is a data-independent
scoring function by the time it is conformalized on $\mathcal D_{\mathrm{calib}}$, so
Theorem~\ref{thm:coverage} applies directly: $P(y_{\mathrm{test}}\in\hat C(\bm x_{\mathrm{test}}))
\ge 1-\alpha$ exactly, independent of whether $f_{\mathrm{mix}}$ or $f_c$ is well specified.
\end{proposition}

\paragraph{Choosing the split.} A single $50/50$ tune/calib split was tried and found empirically
to \emph{lose} to plain MoG-HPD outright on this repository's data: halving the calibration fold
starves $\hat t$'s quantile estimate of data faster than the tuning gain compensates. Two changes
fix this:
\begin{itemize}[leftmargin=1.5em]
\item \textbf{Small tune fraction.} $\texttt{tune\_frac}=0.2$ by default, not $0.5$ -- the
$(c,\theta)$ grid is low-dimensional and discrete, so it needs far less data to rank candidates
than the continuous $(1-\alpha)$ quantile needs to be estimated precisely.
\item \textbf{Repeated splitting with majority vote.} This mirrors the cross-conformal-style
resampling used elsewhere for tuned conformal methods \cite{gibbs2021adaptive}. A single split's
chosen $(c^\star,\theta^\star)$
is noticeably split-dependent (different random tune/calib partitions can pick different grid
cells). Draw $n_{\mathrm{repeats}}$ (default $10$) independent splits; take $(c^\star,\theta^\star)$
as the \emph{mode} (majority vote) of the per-split arg-mins -- appropriate for a discrete grid,
where an arithmetic mean of grid indices would not even necessarily land on a grid point -- and
take the final threshold as the \emph{median} of the per-split calib-fold quantiles at that now-fixed
$(c^\star,\theta^\star)$:
\begin{equation}
\hat t = \operatorname*{median}_{r=1,\dots,n_{\mathrm{repeats}}}\
\operatorname{Quantile}_{q_{\mathrm{level}}}\Big(\big\{S_i(c^\star,\theta^\star)\big\}_{i\in\mathcal
D_{\mathrm{calib}}^{(r)}}\Big).
\label{eq:medianquantile}
\end{equation}
\end{itemize}
Because $(1,0)$ is always on the grid, the tune-fold objective's arg-min can never be worse than
evaluating it at $(1,0)$ alone: STA-HPD's own tuning search can never underperform plain MoG-HPD
on the fold used to pick $(c^\star,\theta^\star)$, and empirically this floor also holds on the
held-out test data across every dataset tried (Section~\ref{sec:results}).

\begin{algorithm}[t]
\caption{STA-HPD fitting}
\label{alg:stahpd}
\begin{algorithmic}[1]
\State \textbf{Input:} calibration set $\mathcal D_{\mathrm{cal}}$, grids $\mathcal C,\Theta$,
$\alpha$, $n_{\mathrm{repeats}}$
\For{$r=1$ to $n_{\mathrm{repeats}}$}
  \State split $\mathcal D_{\mathrm{cal}}$ into $\mathcal D_{\mathrm{tune}}^{(r)}, \mathcal
D_{\mathrm{calib}}^{(r)}$
  \For{$(c,\theta)\in\mathcal C\times\Theta$}
    \State fit $\hat t^{(r)}(c,\theta)$ and measure mean region width on $\mathcal
D_{\mathrm{tune}}^{(r)}$
  \EndFor
  \State $(c^{(r)},\theta^{(r)}) \gets \arg\min$ tune-fold width
\EndFor
\State $(c^\star,\theta^\star) \gets$ majority vote over $\{(c^{(r)},\theta^{(r)})\}_{r=1}^{n_{\mathrm{repeats}}}$
\For{$r=1$ to $n_{\mathrm{repeats}}$}
  \State $\hat t^{(r)} \gets$ quantile of $\{S_i(c^\star,\theta^\star)\}_{i\in\mathcal
D_{\mathrm{calib}}^{(r)}}$
\EndFor
\State $\hat t \gets \operatorname{median}_r \hat t^{(r)}$
\State \textbf{Output:} $(c^\star,\theta^\star,\hat t)$; predict via
Eq.~\eqref{eq:staregion} using Section~\ref{sec:numerics}'s grid-scan-plus-Brent root finder
\end{algorithmic}
\end{algorithm}

\subsection{STA-HPD's results}
\label{sec:results}

Table~\ref{tab:stahpd-results} reports STA-HPD against plain MoG-HPD on four tabular regression
benchmarks with a $K=3$-expert model trained by mixture negative log-likelihood, $1-\alpha=90\%$
nominal coverage, $\texttt{tune\_frac}=0.2$, $n_{\mathrm{repeats}}=10$, over the grid
$\mathcal C=\{0.2,0.4,\dots,2.0\}$, $\Theta=\{-1.0,-0.75,\dots,3.0\}$ (numbers reproduced from this
repository's logged calibration runs). Realized coverage across all cells stayed within
$[88.6\%,91.5\%]$ of the $90\%$ target, so the width comparison below is not confounded by any
method quietly under-covering to buy a narrower interval.

\begin{table}[h]
\centering
\caption{Average interval width, MoG-HPD vs.\ STA-HPD, at matched $\sim\!90\%$ coverage.}
\label{tab:stahpd-results}
\begin{tabular}{@{}lrrrl@{}}
\toprule
\textbf{Dataset} & \textbf{MoG-HPD} & \textbf{STA-HPD} & $\Delta$ & $(c^\star,\theta^\star)$ \\
\midrule
Synthetic & 1.0490 & \textbf{0.9925} & $-5.4\%$ & $(1.3, 0.25)$ \\
Bike & 0.7940 & \textbf{0.7729} & $-2.7\%$ & $(1.0, 0.75)$ \\
Temperature & 0.9817 & \textbf{0.9673} & $-1.5\%$ & $(1.3, 0.25)$ \\
Superconductivity & 0.7964 & \textbf{0.7816} & $-1.9\%$ & $(0.8, 0.25)$ \\
\bottomrule
\end{tabular}
\end{table}

The reduction is positive on all four datasets (and no dataset's search collapses back onto
$(c^\star,\theta^\star)=(1,0)$), with the largest gain on Synthetic -- the one dataset here whose
true conditional law is a known, genuine mixture, exactly the regime Section~\ref{sec:hpd}
motivated highest-density regions for in the first place.

% ============================================================================
\section{SETA-HPD: also correcting the between-expert spread}
\label{sec:seta}

\subsection{The gap STA-HPD still leaves open}

Look back at Eq.~\eqref{eq:sigmatotc}: $\sigma_{\mathrm{tot}}^2(\bm x;c)$ already contains
\emph{both} terms of the aleatoric/epistemic decomposition (Eq.~\ref{eq:lotv}) -- but $c$ (via
Eq.~\ref{eq:sigmashape}) only ever reshapes the first, the within-expert term $\sum_k
w_k\tilde\sigma_k^2$. The second term, $\sum_k w_k(\mu_k-\hat y)^2$ -- \emph{how far apart the gate
places the expert means} -- is left in Eq.~\eqref{eq:fc} and~\eqref{eq:staregion} exactly as the
base model reported it, however systematically over- or under-separated that spread actually is.

This matters concretely for a Mixture-of-Experts model, for a specific, mechanical reason: an MoE
trained with a load-balancing penalty (needed to stop the gate from collapsing all its weight
onto a single expert during training) is under pressure that can, as a side effect, push
well-separated experts \emph{further} apart than the true conditional law actually warrants. That
is precisely the kind of artificially deep density valley that Section~\ref{sec:hpd}'s discussion
of Example~\ref{ex:running} showed is expensive for a level-set method to cover -- except now the
valley is partly an artifact of training pressure, not a real feature of the outcome's
distribution. \textbf{SETA-HPD} (Shape--Epistemic--Threshold Adaptive HPD) adds the missing third
knob to correct exactly this.

\subsection{The epistemic mean-shift}

Introduce a second exponent $\rho \ge 0$, acting on the between-expert term only, through a shift
of the expert means themselves (not their scales):
\begin{equation}
\tilde\mu_k(\bm x;\rho) = \hat y(\bm x) + \sqrt\rho\,\big(\mu_k(\bm x) - \hat y(\bm x)\big).
\label{eq:etamu}
\end{equation}
Each expert's mean is pulled toward (or pushed away from) the aggregate prediction $\hat y(\bm x)$,
by a factor of $\sqrt\rho$ applied to its original distance from $\hat y(\bm x)$.

\begin{lemma}[What the mean-shift does, exactly]
\label{lem:etashift}
For every $\bm x$: $\sum_k w_k(\bm x)\,\tilde\mu_k(\bm x;\rho) = \hat y(\bm x)$ (the aggregate
prediction is exactly unchanged, for \emph{any} $\rho$), and $\sum_k w_k(\bm x)\big(\tilde\mu_k(\bm
x;\rho) - \hat y(\bm x)\big)^2 = \rho\,\sum_k w_k(\bm x)\big(\mu_k(\bm x)-\hat y(\bm x)\big)^2$
(the epistemic variance is scaled by exactly $\rho$). Consequently $\tilde\mu_k(\bm x;1)=\mu_k(\bm
x)$ (identity) and $\tilde\mu_k(\bm x;0)=\hat y(\bm x)$ for every $k$ (every expert collapses onto
the single aggregate prediction, i.e.\ the epistemic spread is switched off entirely, mirroring
what $c=0$ did to the aleatoric term in Lemma~\ref{lem:cdoes}).
\end{lemma}
\begin{proof}
Substitute Eq.~\eqref{eq:etamu} directly. For the mean: $\sum_k w_k\tilde\mu_k = \sum_k w_k\hat y +
\sqrt\rho\sum_k w_k(\mu_k-\hat y) = \hat y + \sqrt\rho\,(\hat y-\hat y) = \hat y$, using $\sum_k
w_k=1$ and $\hat y=\sum_k w_k\mu_k$ (Eq.~\ref{eq:yhatdef}). For the variance:
$\sum_k w_k(\tilde\mu_k-\hat y)^2 = \sum_k w_k\big(\sqrt\rho(\mu_k-\hat y)\big)^2 = \rho\sum_k
w_k(\mu_k-\hat y)^2$ directly. Substituting $\rho=1,0$ gives the two boundary identities.
\end{proof}

The first half of Lemma~\ref{lem:etashift} is the important structural guarantee: because the
shift is built to be mean-\emph{preserving}, applying it never changes the model's point
prediction or, downstream, the target coverage level -- it only redistributes where, around that
fixed point, the mixture's mass sits. Combined with Lemma~\ref{lem:cdoes}, $c$ and $\rho$ act on
\emph{disjoint} pieces of the variance decomposition (Eq.~\ref{eq:lotv}): $c$ reshapes only
$\sigma_a^2$, $\rho$ reshapes only $\sigma_e^2$, and neither touches $\hat y(\bm x)$.

\begin{example}[Example~\ref{ex:running}, concluded]
Recall $\hat y=3.6$, $\mu_1=2$, $\mu_2=6$. At $\rho=0.25$, $\sqrt\rho=0.5$:
\[
\tilde\mu_1(0.25) = 3.6 + 0.5\,(2-3.6) = 3.6 - 0.8 = 2.8, \qquad
\tilde\mu_2(0.25) = 3.6 + 0.5\,(6-3.6) = 3.6 + 1.2 = 4.8.
\]
The gap between the two expert means shrinks from $6-2=4$ to $4.8-2.8=2$ -- exactly halved, since
$\sqrt{0.25}=0.5$ -- while their weighted average is still exactly $0.6\times2.8+0.4\times4.8 =
1.68+1.92=3.6=\hat y$, unchanged. This directly fills in part of the density valley between the
two modes that Section~\ref{sec:hpd} identified as the source of wasted width: at $\rho=1$ the
valley sits at the true $\hat y=3.6$ between well-separated modes at $2$ and $6$; at $\rho=0.25$
the same valley sits between modes at $2.8$ and $4.8$, a much narrower gap for a level-set method
to have to cover or split across.
\end{example}

\subsection{The SETA-HPD family}

Combine the reshaped scale $\tilde\sigma_k(\bm x;c)$ (Eq.~\ref{eq:sigmashape}, unchanged from
STA-HPD) with the shifted mean $\tilde\mu_k(\bm x;\rho)$ (Eq.~\ref{eq:etamu}, new) into a jointly
reshaped density and total scale,
\begin{align}
f_{c,\rho}(y\mid\bm x) &= \sum_{k=1}^K w_k(\bm x)\,\mathcal N\big(y;\ \tilde\mu_k(\bm x;\rho),\
\tilde\sigma_k^2(\bm x;c)\big), \label{eq:fcrho}\\
\sigma_{\mathrm{tot}}^2(\bm x;c,\rho) &= \sum_{k=1}^K w_k(\bm x)\,\tilde\sigma_k^2(\bm x;c) \;+\;
\rho\sum_{k=1}^K w_k(\bm x)\,\big(\mu_k(\bm x)-\hat y(\bm x)\big)^2, \label{eq:sigmatotrho}
\end{align}
where the second line follows directly from Lemma~\ref{lem:etashift}'s exact-$\rho$-scaling
identity (there is no need to recompute the epistemic term from the shifted means separately -- it
is, by construction, exactly $\rho$ times the original). The score, threshold, and induced region
are the three-exponent analogues of Eqs.~\eqref{eq:stascore}--\eqref{eq:staregion}, with every
occurrence of $\theta$'s field now built from $\sigma_{\mathrm{tot}}(\bm x;c,\rho)$ instead of
$\sigma_{\mathrm{tot}}(\bm x;c)$:
\begin{align}
S_i(c,\rho,\theta) &= -\log f_{c,\rho}(y_i\mid\bm x_i) - \theta\log\sigma_{\mathrm{tot}}(\bm
x_i;c,\rho), \label{eq:setascore}\\
\hat C(\bm x) &= \Big\{y : \log f_{c,\rho}(y\mid\bm x) \ge -\hat t(c,\rho,\theta) -
\theta\log\sigma_{\mathrm{tot}}(\bm x;c,\rho)\Big\}. \label{eq:setaregion}
\end{align}

\subsection{Boundary recovery}

\begin{proposition}[Boundary recovery, three exponents]
\label{prop:setarecovery}
$(c,\rho,\theta)=(1,1,0)$ recovers MoG-HPD exactly. $(c,\rho,\theta)=(c,1,\theta)$, for
\emph{any} $(c,\theta)$, recovers STA-HPD exactly -- so STA-HPD is not a single point of
SETA-HPD's family but an entire sub-family (the whole $\rho=1$ slice). $(c,\rho,\theta)=(1,\rho,0)$,
for any $\rho$, is a new intermediate calibrator that scales only the epistemic spread, leaving the
aleatoric term and the threshold untouched. $(c,\rho,\theta)=(1,1,1)$ recovers CP-VS exactly when
$K=1$ (where the epistemic term is identically zero, so $\rho$ has no effect at all).
\end{proposition}
\begin{proof}
Each case is direct substitution into Eqs.~\eqref{eq:etamu}--\eqref{eq:setaregion}, using
Lemma~\ref{lem:etashift} ($\tilde\mu_k(\bm x;1)=\mu_k(\bm x)$) and Lemma~\ref{lem:cdoes}
($\tilde\sigma_k(\bm x;1)=\sigma_k(\bm x)$), exactly as in the proof of
Proposition~\ref{prop:starecovery}.
\end{proof}

Table~\ref{tab:cheatsheet} in Section~\ref{sec:summary} collects every method introduced in this
document as a single row of the $(c,\rho,\theta)$ grid, for quick reference.

Because $\rho=1$ is always kept on SETA-HPD's search grid (by construction: even if a user passes
a custom grid that omits it, $1.0$ is re-inserted), and $(1,0)$ is always on STA-HPD's own
$(c,\theta)$ grid, the arg-min of SETA-HPD's tuning objective over the full 3-D grid can never
exceed the arg-min restricted to the $\rho=1$ slice -- which is exactly STA-HPD's own tuning
objective. So SETA-HPD's tuning-fold performance can never be worse than STA-HPD's, which in turn
can never be worse than MoG-HPD's: the three methods form a strict, nested floor. Validity is
completely unaffected by adding the third exponent: Proposition~\ref{prop:stavalidity}'s argument
applies verbatim with $(c,\rho,\theta)$ selected on $\mathcal D_{\mathrm{tune}}$ in place of
$(c,\theta)$, since nothing about the exchangeability argument in Theorem~\ref{thm:coverage}
depends on how many exponents the frozen, label-independent scoring function happens to have.

\subsection{Fitting}

The fitting procedure is Algorithm~\ref{alg:stahpd} verbatim, with the grid searched at each split
widened from $\mathcal C\times\Theta$ to $\mathcal C\times\mathcal R\times\Theta$, where $\mathcal
R = \{0,0.25,0.5,0.75,1.0,1.5,2.0\}$ by default ($1.0$ always included, for the boundary-recovery
argument above to hold). One implementation detail keeps the added cost of the third exponent
small: $f_{c,\rho}$ does not depend on $\theta$ at all (Eq.~\ref{eq:fcrho}), so for each $(c,\rho)$
pair the reshaped density is evaluated on the tune fold exactly once, and the entire $\theta$ grid
is then swept against it using only cheap quantile and threshold-count operations -- the expensive
part (evaluating a $K$-component mixture density on a fine grid) does not get repeated once per
$\theta$, keeping the added cost roughly proportional to $|\mathcal R|$ rather than $|\mathcal
R|\times|\Theta|$.

\subsection{SETA-HPD's results, and when the third knob actually helps}
\label{sec:setaresults}

Table~\ref{tab:setaresults} extends Table~\ref{tab:stahpd-results} with SETA-HPD, same checkpoints,
$1-\alpha=90\%$, matched tune/calib splits so STA-HPD and SETA-HPD are compared on identical folds.

\begin{table}[h]
\centering
\caption{MoG-HPD vs.\ STA-HPD vs.\ SETA-HPD at matched coverage. $\rho^\star=1$ means the search
left the epistemic term untouched (SETA-HPD collapsed back onto its own STA-HPD sub-family).}
\label{tab:setaresults}
\begin{tabular}{@{}lrrrrr@{}}
\toprule
\textbf{Dataset} & \textbf{MoG-HPD} & \textbf{STA-HPD} & \textbf{SETA-HPD} & $\Delta$ vs.\ STA-HPD
& $\rho^\star$ \\
\midrule
Synthetic & 1.0490 & 0.9925 & 0.9925 & $0.0\%$ & $1.00$ \\
Bike & 0.7940 & 0.7729 & \textbf{0.7399} & $-4.3\%$ & $0.25$ \\
Temperature & 0.9817 & 0.9673 & 0.9673 & $0.0\%$ & $1.00$ \\
Superconductivity & 0.7964 & 0.7816 & 0.7849 & $+0.4\%$ & $0.75$ \\
\bottomrule
\end{tabular}
\end{table}

The change relative to STA-HPD is never negative by more than measurement noise, consistent with
the nested-floor argument above, but the practical picture is dataset-dependent rather than a
uniform win:
\begin{itemize}[leftmargin=1.5em]
\item On \textbf{Synthetic} and \textbf{Temperature}, the search selects $\rho^\star=1$ and stays
there across repeated resamples: STA-HPD's own $(c,\theta)$ already capture whatever is correctable
about these two datasets, and there is no further, separately-correctable epistemic miscalibration
left for $\rho$ to fix.
\item On \textbf{Bike}, $\rho^\star=0.25$ is selected stably and yields the largest gain in the
sweep, $-4.3\%$ width relative to STA-HPD, while also eliminating disjoint multi-segment prediction
sets entirely (from a small but nonzero fraction of test points down to none): compressing the two
experts toward $\hat y$ removes the density valley that was fragmenting some predictions into two
separate intervals, exactly the mechanism illustrated in the running example above.
\item On \textbf{Superconductivity}, $\rho^\star=0.75$ is selected stably; mean width is
essentially flat, but the fraction of fragmented (disjoint, multi-interval) prediction sets drops
substantially -- the same coverage budget is being spent as fewer, more usable single intervals
rather than as raw width reduction, a benefit not visible in the width column alone.
\end{itemize}
The general lesson: $\rho$ is a genuinely separate degree of freedom from $(c,\theta)$ (it acts on
a different, disjoint term of the variance decomposition, per Lemma~\ref{lem:etashift}), but
whether that specific degree of freedom is the bottleneck on a given dataset -- as opposed to the
aleatoric shape or the threshold level, which $(c,\theta)$ already correct -- is itself a
data-dependent question, which is exactly what the tuning search in Section~\ref{sec:stafit}
(extended to the 3-D grid) is built to answer safely, i.e.\ without ever risking a worse result
than not using $\rho$ at all.

Section~\ref{sec:stafit}'s tuning search picks every exponent -- $(c,\theta)$ for STA-HPD,
$(c,\rho,\theta)$ for SETA-HPD -- by the same objective: minimize mean prediction-set width on a
tune fold $\mathcal D_{\mathrm{tune}}$. That width is a noisy, resampled estimate of a population
quantity, and giving the search one more dimension to explore (as SETA-HPD's $\rho$ does) gives it
more room to fit that noise rather than genuine signal, so a $(c^\star,\rho^\star,\theta^\star)$
that merely got lucky on one tune fold can occasionally generalize worse, on held-out data, than
the plainer STA-HPD it was meant to improve on. The next two sections pick up this thread from two
different directions: Section~\ref{sec:ease} (EASE-HPD) widens the same width-driven search by one
further axis, while Section~\ref{sec:meld} (MELD-HPD) instead replaces the width objective itself,
for the two exponents that actually define the surrogate density, with a different, smoother
criterion.

% ============================================================================
\section{EASE-HPD: Letting the Threshold Respond to \emph{Composition}, Not Just Total Scale}
\label{sec:ease}

\subsection{The gap SETA-HPD still leaves open}

Look once more at $\sigma_{\mathrm{tot}}^2(\bm x;c,\rho)$ from Eq.~\eqref{eq:sigmatotrho}. It is a
\emph{sum} of an aleatoric piece and an epistemic piece; write them out separately,
\begin{equation}
A(\bm x;c) = \sum_{k=1}^K w_k(\bm x)\,\tilde\sigma_k^2(\bm x;c), \qquad
E(\bm x;\rho) = \sum_{k=1}^K w_k(\bm x)\,\big(\tilde\mu_k(\bm x;\rho) - \hat y(\bm x)\big)^2,
\label{eq:aeterms}
\end{equation}
so that $\sigma_{\mathrm{tot}}^2(\bm x;c,\rho) = A(\bm x;c) + E(\bm x;\rho)$. The threshold-field
exponent $\theta$ (Section~\ref{sec:stahpd}) multiplies $\log\sigma_{\mathrm{tot}}$ -- it only ever
sees this \emph{sum}. Two test points can have exactly the same total $A+E$ while having
completely different \emph{compositions}: one point's uncertainty might be almost entirely
within-expert noise ($E\approx0$), another's almost entirely expert disagreement ($A\approx0$).
$\theta$, as defined so far, treats these two points identically, because it never looks past the
sum. That is a real gap, for a specific, mechanical reason particular to this MoE setting: as
already discussed in Section~\ref{sec:seta}, the epistemic term $E(\bm x;\rho)$ is shaped in part
by the load-balancing pressure used during training (an artifact of how the gate is regularized),
while the aleatoric term $A(\bm x;c)$ is fit directly against the data by each expert's own
negative log-likelihood and is typically much better behaved. If the threshold could tell \emph{how
much} of a point's total scale comes from the (less trustworthy) epistemic piece versus the (more
trustworthy) aleatoric piece, it could tilt accordingly. \textbf{EASE-HPD} (Epistemic-Aleatoric
Share-adaptive threshold HPD) adds exactly this fourth knob.

\subsection{The epistemic share}

Define the (post-reshaping) \emph{epistemic share} of a point's total predictive variance,
\begin{equation}
s(\bm x;c,\rho) = \frac{E(\bm x;\rho)}{A(\bm x;c) + E(\bm x;\rho)} =
\frac{E(\bm x;\rho)}{\sigma_{\mathrm{tot}}^2(\bm x;c,\rho)} \ \in\ [0,1].
\label{eq:share}
\end{equation}
$s(\bm x;c,\rho)=0$ means ``all of this point's uncertainty is within-expert noise''; $s(\bm
x;c,\rho)=1$ means ``all of it is expert disagreement''; a value like $0.35$ means roughly a third
of the reported variance at $\bm x$ comes from experts disagreeing rather than from any one
expert's own reported spread.

\begin{lemma}[The share is always well defined, and collapses exactly when the epistemic term
does]
\label{lem:share}
For every $\bm x,c,\rho$: $A(\bm x;c)\ge0$ and $E(\bm x;\rho)\ge0$ (both are weighted sums of
squares with nonnegative weights $w_k$), so $s(\bm x;c,\rho)$ is a well-defined ratio in $[0,1]$
whenever $\sigma_{\mathrm{tot}}(\bm x;c,\rho)>0$, which always holds since every $\tilde\sigma_k(\bm
x;c)>0$. Moreover $s(\bm x;c,\rho)\equiv0$ for every $\bm x$ whenever $\rho=0$ \emph{or} $K=1$.
\end{lemma}
\begin{proof}
Well-definedness and the $[0,1]$ range follow directly from $A,E\ge0$. At $\rho=0$,
Lemma~\ref{lem:etashift} gives $\tilde\mu_k(\bm x;0)=\hat y(\bm x)$ for every $k$, so
$E(\bm x;0) = \sum_k w_k(\hat y(\bm x)-\hat y(\bm x))^2 = 0$ for every $\bm x$. At $K=1$,
$\hat y(\bm x)=\mu_1(\bm x)$ identically (there is only one expert to average), so
$\tilde\mu_1(\bm x;\rho)-\hat y(\bm x)=0$ for every $\rho$, giving $E(\bm x;\rho)=0$ again. In both
cases $s(\bm x;c,\rho) = 0/A(\bm x;c) = 0$.
\end{proof}

This lemma matters practically, not just formally: it says in advance, precisely, the two
situations in which the new knob introduced below is guaranteed to do nothing -- so that when it is
later observed to do nothing on a particular dataset (Section~\ref{sec:easeexample}), that is a
predicted consequence of the construction, not a bug to debug.

\subsection{The EASE-HPD family}

Freeze the calibration-set mean share at each $(c,\rho)$ cell, $\bar s(c,\rho) = \tfrac1n\sum_i
s(\bm x_i;c,\rho)$ (exactly as $G$ was frozen for $c$ in Eq.~\ref{eq:Gdef}). A fourth exponent
$\phi\in\mathbb R$ tilts the score by a point's own \emph{deviation} of its share from that typical
value,
\begin{align}
S_i(c,\rho,\theta,\phi) &= -\log f_{c,\rho}(y_i\mid\bm x_i) \;-\; \theta\log\sigma_{\mathrm{tot}}(\bm
x_i;c,\rho) \;-\; \phi\big(s(\bm x_i;c,\rho) - \bar s(c,\rho)\big), \label{eq:easescore}\\
\hat C(\bm x) &= \Big\{y : \log f_{c,\rho}(y\mid\bm x) \ge -\hat t(c,\rho,\theta,\phi) -
\theta\log\sigma_{\mathrm{tot}}(\bm x;c,\rho) - \phi\big(s(\bm x;c,\rho)-\bar s(c,\rho)\big)
\Big\}. \label{eq:easeregion}
\end{align}
Centering by $\bar s(c,\rho)$ is what keeps $\phi$ from simply fighting $\hat t$ over the same
additive constant (without centering, a nonzero average share would let $\phi$ silently shift the
effective threshold for \emph{every} point, not just tilt it based on how a point's composition
differs from typical); boundedness of $s$ in $[0,1]$ (Lemma~\ref{lem:share}) means this term can
never blow up the way an uncentered ratio $E/A$ could as $A\to0$.

\begin{proposition}[Boundary recovery, four exponents]
\label{prop:easerecovery}
$(c,\rho,\theta,\phi)=(1,1,0,0)$ recovers MoG-HPD exactly. $(c,\rho,\theta,0)$, for \emph{any}
$(c,\rho,\theta)$, recovers SETA-HPD exactly -- so SETA-HPD, and by
Proposition~\ref{prop:setarecovery} STA-HPD and MoG-HPD with it, are entire sub-families of
EASE-HPD, not merely isolated points.
\end{proposition}
\begin{proof}
At $\phi=0$ the term $-\phi(s(\bm x;c,\rho)-\bar s(c,\rho))$ vanishes identically in
Eqs.~\eqref{eq:easescore}--\eqref{eq:easeregion}, regardless of the value of $s$, leaving exactly
SETA-HPD's score and region (Eqs.~\ref{eq:setascore}--\ref{eq:setaregion}); the case
$(1,1,0,0)$ then follows from Proposition~\ref{prop:setarecovery}'s own $(1,1,0)$ case.
\end{proof}

\begin{remark}[$\phi$ is provably inert exactly where Lemma~\ref{lem:share} predicts]
\label{rem:phiinert}
Whenever the tuned $\rho^\star=0$, or whenever $K=1$, Lemma~\ref{lem:share} gives
$s(\bm x;c,\rho)\equiv0$ for every calibration point, hence $\bar s(c,\rho)=0$ too, so
$s(\bm x;c,\rho)-\bar s(c,\rho)\equiv0$ identically: the $\phi$ term contributes exactly zero to
Eqs.~\eqref{eq:easescore}--\eqref{eq:easeregion} for \emph{every} value of $\phi$. In that regime
$\phi$ cannot be identified from the tuning objective at all (every candidate $\phi$ scores
identically), and whichever value the tie-break happens to pick has literally no effect on the
output.
\end{remark}

Validity carries over unchanged: Proposition~\ref{prop:stavalidity}'s argument applies verbatim
with $(c,\rho,\theta,\phi)$ selected on $\mathcal D_{\mathrm{tune}}$ in place of $(c,\theta)$, since
$\phi(s(\bm x;c,\rho)-\bar s(c,\rho))$ is, for any fixed $(c,\rho,\theta,\phi)$, simply another
data-independent function of $\bm x$ -- adding it changes nothing about the exchangeability
argument behind Theorem~\ref{thm:coverage}.

\paragraph{Fitting.} The search widens from $\mathcal C\times\mathcal R\times\Theta$ to $\mathcal
C\times\mathcal R\times\Theta\times\Phi$, with $\Phi=\{-1.0,-0.5,0,0.5,1.0\}$ by default ($0\in\Phi$
always kept, for the same boundary-recovery reason as every other knob in this document). Since
neither $f_{c,\rho}$ nor $s(\cdot;c,\rho)$ depends on $\theta$ or $\phi$, both are computed once per
$(c,\rho)$ pair and the whole $(\theta,\phi)$ sub-grid is then swept with cheap quantile and
threshold-count operations, mirroring the same factorization trick used for SETA-HPD's grid
(Section~\ref{sec:seta}). Because $\rho=1$ and $\phi=0$ are always on their respective grids, the
tune-fold objective can never be worse at the arg-min than at $(c,\rho,\theta,0)$ for any
$(c,\rho,\theta)\in\mathcal C\times\mathcal R\times\Theta$, which in turn can never be worse than
$(1,1,0,0)$: EASE-HPD's tuning objective can never exceed SETA-HPD's, which can never exceed
STA-HPD's or plain MoG-HPD's -- the same nested floor as Section~\ref{sec:seta}, now with one more
link in the chain.

\subsection{A worked example, including the collapse case}
\label{sec:easeexample}

On the single-expert Synthetic checkpoint used elsewhere in this document (seed $4022$,
$1-\alpha=90\%$), the fitting search selects $(c^\star,\rho^\star,\theta^\star,\phi^\star) =
(0.8,\,0,\,1.5,\,-1.0)$, with the tune-fold width identical across \emph{every} candidate $\phi$
(the logged $\phi$-marginal reads $1.4402$ at all five grid values). This is
Remark~\ref{rem:phiinert}'s collapse case, playing out on real data: the search independently lands
on $\rho^\star=0$ (the same choice SETA-HPD makes on its own on this dataset,
Table~\ref{tab:setaresults}), which by Lemma~\ref{lem:share} forces $s(\bm x;c,0)\equiv0$ and $\bar
s=0$ -- so $\phi$ is provably unable to move either the score or the region here, regardless of
which value it happens to be assigned. EASE-HPD's realized coverage and width on this checkpoint,
$0.8933$ and $1.3473$, are accordingly identical to STA-HPD's and SETA-HPD's to four decimal
places (Table~\ref{tab:allfive} below) -- an empirical confirmation, not merely a theoretical
prediction, of Remark~\ref{rem:phiinert}.

% ============================================================================
\section{MELD-HPD: Selecting Shape by a Different Objective Entirely}
\label{sec:meld}

\subsection{Why picking $(c,\rho)$ by width can misfire}

Every method up to this point -- STA-HPD, SETA-HPD, EASE-HPD -- selects its exponents by the same
recipe: minimize an estimate $\bar W$ of the tune-fold's mean region width
(Section~\ref{sec:stafit}), computed on a \emph{subsample} of the tune fold using the fast
fixed-grid approximation of Section~\ref{sec:numerics}. This estimate is a Monte-Carlo
approximation of a population quantity, not the population quantity itself, and it is built from a
capped, subsampled set of points for speed. Widening the grid that gets searched (as each of
$c\to(c,\theta)\to(c,\rho,\theta)\to(c,\rho,\theta,\phi)$ did) gives the arg-min more room to fit
sampling noise \emph{in that specific noisy estimator}, not necessarily more room to reduce the
population width it is approximating. In practice this shows up exactly as this risk predicts:
SETA-HPD is documented (Section~\ref{sec:setaresults}, and this repository's own diagnostic logs)
to occasionally come out \emph{wider than STA-HPD on held-out test data}, on specific
(architecture, dataset) configurations, despite the tune-fold floor $\bar W(c^\star,\rho^\star,
\theta^\star)\le\bar W(1,1,0)$ holding exactly, by construction, on the very fold that chose those
parameters. The floor Section~\ref{sec:seta} establishes is real -- but it is a floor on the
\emph{estimator}, evaluated where it was fit, not a guarantee about the \emph{population} quantity
it approximates, evaluated somewhere new. \textbf{MELD-HPD} (Maximum-likelihood Epistemic-aLeatoric
Decomposition) responds by selecting the two \emph{density-shaping} exponents $(c,\rho)$ --
the only two that change $f_{c,\rho}$ itself, as opposed to $\theta$, which only ever tilts a
threshold around an already-fixed density -- using a completely different, smoother objective,
evaluated on the whole tune fold rather than a width subsample.

\subsection{The log-score objective, and what it is secretly minimizing}

Define the mean calibration \emph{log-density} on the tune fold,
\begin{equation}
L(c,\rho) = \frac{1}{|\mathcal D_{\mathrm{tune}}|} \sum_{i\in\mathcal D_{\mathrm{tune}}} \log
f_{c,\rho}(y_i\mid\bm x_i).
\label{eq:logscore}
\end{equation}
This is simply the average log-likelihood the reshaped density $f_{c,\rho}$ assigns to the actual
observed outcomes -- a completely standard model-fit criterion, with no reference to conformal
prediction, width, or coverage at all. The next result says precisely what maximizing it targets.

\begin{proposition}[The log score targets the KL-closest member of the family]
\label{prop:logscore-kl}
Let $p(\cdot\mid\bm x)$ denote the true (unknown) conditional density of $Y\mid\bm x$. Then
\begin{equation}
\mathbb E\big[\log f_{c,\rho}(Y\mid X)\big] = -\,\mathbb E_X\big[H(p(\cdot\mid X))\big] \;-\;
\mathbb E_X\big[\mathrm{KL}\big(p(\cdot\mid X)\,\big\|\,f_{c,\rho}(\cdot\mid X)\big)\big],
\label{eq:klidentity}
\end{equation}
where $H$ is differential entropy and $\mathrm{KL}$ is Kullback--Leibler divergence -- a standard
measure of how different two densities are, always $\ge0$, equal to $0$ exactly when the two
densities agree almost everywhere. The first term on the right does not depend on $(c,\rho)$ at
all, so maximizing the left side over $(c,\rho)$ is \emph{exactly} minimizing $\mathbb
E_X[\mathrm{KL}(p(\cdot\mid X)\,\|\,f_{c,\rho}(\cdot\mid X))]$: the population log score is
maximized precisely at whichever $(c,\rho)$ makes the reshaped density closest, on average in KL
divergence, to the true conditional law.
\end{proposition}
\begin{proof}
For a fixed $\bm x$, split $\int p(y\mid\bm x)\log f_{c,\rho}(y\mid\bm x)\,dy$ by adding and
subtracting $\int p(y\mid\bm x)\log p(y\mid\bm x)\,dy$:
\[
\int p(y\mid\bm x)\log f_{c,\rho}(y\mid\bm x)\,dy = \int p(y\mid\bm x)\log p(y\mid\bm x)\,dy -
\int p(y\mid\bm x)\log\frac{p(y\mid\bm x)}{f_{c,\rho}(y\mid\bm x)}\,dy = -H(p(\cdot\mid\bm x)) -
\mathrm{KL}\big(p(\cdot\mid\bm x)\,\|\,f_{c,\rho}(\cdot\mid\bm x)\big),
\]
the standard cross-entropy/entropy/KL decomposition. Taking $\mathbb E_X$ of both sides gives
Eq.~\eqref{eq:klidentity}. Since $\mathrm{KL}\ge0$ and $\mathbb E_X[H(p(\cdot\mid X))]$ is a fixed
constant of the (unknown, but fixed) true law, maximizing the left-hand side over $(c,\rho)$ is
exactly minimizing the KL term.
\end{proof}

This is the population-level reason the log score is called a \emph{strictly proper} scoring rule
\cite{gneiting2007strictly}: no other member of the family can achieve a higher expected score than
the KL-closest one, so maximizing the sample version $L(c,\rho)$ (Eq.~\ref{eq:logscore}) is a
sensible, well-founded way to estimate the population-optimal $(c,\rho)$, using the entire tune
fold rather than a width subsample, with no dependence on any grid-resolution approximation of a
region's measure (unlike $\bar W$, which additionally discretizes the region on a fixed grid).

\begin{remark}[A proxy for width optimality, not a guarantee of it]
\label{rem:logscoreproxy}
Theorem~\ref{thm:optimality}'s optimality condition (Eq.~\ref{eq:optcond}) is a statement about the
true density specifically \emph{at the region's boundary}; it is not, in general, the same
condition as minimizing the \emph{global} average discrepancy $\mathbb E_X[\mathrm{KL}(p(\cdot\mid
X)\|f_{c,\rho}(\cdot\mid X))]$ that Proposition~\ref{prop:logscore-kl} shows the log score targets
-- a density can be globally close in KL divergence while still being systematically wrong exactly
at the region's boundary, or the reverse. Selecting $(c,\rho)$ by the log score is therefore a
well-motivated \emph{proxy} for solving Eq.~\eqref{eq:optcond}, not a construction that provably
achieves it. This is the precise, technical sense in which MELD-HPD's shape selection is
theoretically well grounded but, unlike Proposition~\ref{prop:stavalidity}'s coverage guarantee,
carries no exact guarantee of narrower width -- and the worked comparison below
(Section~\ref{sec:meldexample}) shows this gap is not just a theoretical possibility.
\end{remark}

\subsection{$\theta$ still needs its own objective, and a diagnostic for $\rho$}

$\theta$ never enters $f_{c,\rho}$ (Eq.~\ref{eq:fcrho}) -- it only tilts the threshold around an
already-fixed density -- so $L(c,\rho)$ is constant in $\theta$ and cannot be used to choose it.
$\theta^\star$ is therefore still selected exactly as in STA-HPD/SETA-HPD, by $\arg\min$ tune-fold
width $\bar W(c^\star,\rho^\star,\theta)$, but now at the already-fixed $(c^\star,\rho^\star)$ --
its own proper one-dimensional objective, and a much safer search than a joint $(c,\rho,\theta)$
width arg-min, since only one axis is left for the noisy width estimator to overfit.

As a purely interpretive diagnostic (not part of the selection rule itself), MELD-HPD also reports
a closed-form estimate of how much epistemic spread the residuals actually support:

\begin{lemma}[Least-squares epistemic-reliability estimate]
\label{lem:closedrho}
Fix $c$, write $A_i=A(\bm x_i;c)$, $E_i = \sum_k w_k(\bm x_i)(\mu_k(\bm x_i)-\hat y(\bm x_i))^2$
(the \emph{raw}, un-reshaped, $\rho=1$ epistemic term), and $R_i^2=(y_i-\hat y(\bm x_i))^2$ (the
squared residual). The value of $\rho$ minimizing $J(\rho) = \sum_{i=1}^n (A_i+\rho E_i - R_i^2)^2$
is
\begin{equation}
\hat\rho = \frac{\sum_i E_i(R_i^2-A_i)}{\sum_i E_i^2}, \qquad \text{whenever } \sum_i E_i^2>0,
\label{eq:closedrho}
\end{equation}
and the reported diagnostic is $\max(\hat\rho,0)$ (projected onto $\rho\ge0$, since
Eq.~\eqref{eq:etamu}'s $\sqrt\rho$ requires it). $\hat\rho$ is undefined (reported as \texttt{NaN})
when $\sum_i E_i^2=0$, i.e.\ whenever $K=1$, where the epistemic term carries no information at
all for any $\rho$ to be fit against.
\end{lemma}
\begin{proof}
$J$ is a quadratic, convex function of $\rho$ (its second derivative is $2\sum_i E_i^2\ge0$).
Setting $J'(\rho) = 2\sum_i E_i(A_i+\rho E_i-R_i^2)=0$ and solving gives $\rho\sum_i E_i^2 = \sum_i
E_i(R_i^2-A_i)$, i.e.\ Eq.~\eqref{eq:closedrho}; convexity makes this stationary point the global
minimizer whenever $\sum_i E_i^2>0$.
\end{proof}

Equation~\eqref{eq:closedrho} is exactly the ordinary-least-squares fit, treating $E_i$ as the sole
covariate with fixed offset $A_i$, of the reshaped total-variance identity
$\sigma_{\mathrm{tot}}^2(\bm x;c,\rho) = A(\bm x;c) + \rho E(\bm x)$ against the observed squared
residuals -- in plain terms, ``how much of the model's reported epistemic spread the residuals
actually back up,'' at the log-score-chosen $c^\star$. A value near $\hat\rho\approx1$ says the
reported epistemic spread is about right; $\hat\rho\ll1$ says the experts are more separated than
the residuals warrant (over-separation, the failure mode motivating $\rho$ in the first place,
Section~\ref{sec:seta}); $\hat\rho\gg1$ says the reverse.

\subsection{Boundary recovery, validity, and a weaker floor}

The reachable $(c,\rho,\theta)$ configurations are identical to SETA-HPD's own grid $\mathcal
C\times\mathcal R\times\Theta$, and still contain $(1,1,0)$ (MoG-HPD) and $(c,1,\theta)$ for every
$(c,\theta)$ (STA-HPD), by the same substitutions as Proposition~\ref{prop:setarecovery}.

\begin{remark}[Same grid, a weaker guarantee]
\label{rem:meldweakfloor}
Because $(c^\star,\rho^\star)$ is chosen to \emph{maximize} $L(c,\rho)$ rather than \emph{minimize}
$\bar W(c,\rho,\theta)$, and $(1,1)\in\mathcal C\times\mathcal R$, MELD-HPD's tuning objective only
guarantees $L(c^\star,\rho^\star)\ge L(1,1)$ -- a floor on \emph{log density}, not on width. This
does \emph{not} imply $\bar W(c^\star,\rho^\star,\theta^\star)\le\bar W(1,1,0)$: unlike every other
method in this document, MELD-HPD carries no guarantee that its tune-fold width is at most plain
MoG-HPD's or STA-HPD's -- only that $\theta^\star$ is itself width-optimal \emph{given} the
already-fixed, log-score-selected $(c^\star,\rho^\star)$ (since $\theta=0$ is still on the grid
$\theta^\star$ is chosen from).
\end{remark}

Validity is completely unaffected by any of this: Theorem~\ref{thm:coverage}'s guarantee depends
only on $(c,\rho,\theta)$ being fixed by \emph{some} rule that is a function of the tune fold
alone, independent of the calib fold's labels -- it does not depend on \emph{which} objective that
rule optimizes. Since $(c^\star,\rho^\star)$ is chosen from $\mathcal D_{\mathrm{tune}}$ by
Eq.~\eqref{eq:logscore} and $\theta^\star$ from the same fold by $\bar W$, both disjoint from
$\mathcal D_{\mathrm{calib}}$, the resulting score is a frozen, label-independent function of
$(\bm x_i,y_i)$ by the time it is conformalized, exactly as in
Proposition~\ref{prop:stavalidity}.

\paragraph{Fitting.} Two stages replace the single width arg-min of Algorithm~\ref{alg:stahpd}: (1)
$(c^\star,\rho^\star)$ is the majority vote, over the $n_{\mathrm{repeats}}$ splits, of the
per-split $\arg\max_{c,\rho} L(c,\rho)$ (each per-point log-density is floored at $-10^3$ nats
before averaging, so one numerically underflowing point cannot send an entire grid cell to
$-\infty$); (2) at the now-fixed $(c^\star,\rho^\star)$, $\theta^\star$ is the majority vote of the
per-split $\arg\min_\theta \bar W(c^\star,\rho^\star,\theta)$. The final threshold uses the same
median-of-calib-fold-quantiles construction as every other method here
(Eq.~\ref{eq:medianquantile}), unchanged.

\subsection{A worked comparison across all five methods}
\label{sec:meldexample}

On the same checkpoint, split construction, and nominal coverage used in
Section~\ref{sec:easeexample}, Table~\ref{tab:allfive} reports every method introduced in this
document side by side.

\begin{table}[h]
\centering
\caption{All five calibrators on one held-out checkpoint (Synthetic, single-expert probabilistic
head, seed $4022$, $1-\alpha=90\%$), identical tune/calib construction throughout.}
\label{tab:allfive}
\begin{tabular}{@{}lrrl@{}}
\toprule
\textbf{Method} & \textbf{Coverage} & \textbf{Width} & \textbf{Selected parameters} \\
\midrule
MoG-HPD & $0.9087$ & $1.5757$ & --- \\
STA-HPD & $0.8933$ & $1.3473$ & $(c,\theta)=(0.8,1.5)$ \\
SETA-HPD & $0.8933$ & $1.3473$ & $(c,\rho,\theta)=(0.8,0,1.5)$ \\
EASE-HPD & $0.8933$ & $1.3473$ & $(c,\rho,\theta,\phi)=(0.8,0,1.5,-1.0)$ \\
MELD-HPD & $0.8940$ & $1.3708$ & $(c,\rho,\theta)=(1.3,0,0.75)$ \\
\bottomrule
\end{tabular}
\end{table}

STA-HPD, SETA-HPD, and EASE-HPD agree to four decimal places -- exactly what their respective
nested-floor propositions guarantee on this config, once $\rho^\star=0$ and
Remark~\ref{rem:phiinert}'s inertness are accounted for. MELD-HPD is the one visible exception, at
$1.3708$ -- $1.7\%$ \emph{wider} than the other four -- and the logged tuning surfaces show exactly
why. The log-score surface is monotone increasing in $c$ on this config's grid ($L(0.8,0)=-0.4528$,
$L(1,0)=-0.3885$, $L(1.3,0)=-0.3634$, all at $\rho=0$), so MELD-HPD's proper-scoring-rule selector
correctly, faithfully picks the largest available $c$ on the grid, $c^\star=1.3$. But the tune-fold
\emph{width} surface at that same cell, $\bar W(1.3,0,0.75)=1.4705$, is \emph{not} the grid's
minimum -- that minimum, $\bar W(0.8,0,1.5)=1.4402$, sits at the smaller $c=0.8$ that the other four
methods all select instead. This is Remark~\ref{rem:logscoreproxy} realized on real data: the
density that is closest to the truth, by the strictly proper log score, is demonstrably not the
width-narrowest one on this particular checkpoint -- exactly the gap a KL-motivated selector does
not, and by Remark~\ref{rem:meldweakfloor}'s argument structurally cannot, close.

% ============================================================================
\section{Where these ideas come from, and what is actually new}
\label{sec:related}

It is worth being precise about which parts of this document are borrowed and which are new,
because most of the machinery is not new and it would be misleading to present it as if it were.
Reading the sections above, almost every \emph{ingredient} has a clear home in the existing
literature; what this family contributes is a specific way of combining them.

\paragraph{The level-set idea is not ours.} Turning a predictive density into a prediction set by
keeping its highest-density region (Section~\ref{sec:hpd}) is a classical statistical object
\cite{boxtiao1973bayesian,hyndman1996computing}, and conformalizing it -- so the region has a
distribution-free coverage guarantee even when the density is wrong -- is due to Izbicki, Shimizu
and Stern (CD-split / HPD-split) \cite{izbicki2020cd}, with the classification analogue in Sadinle
et al.\ \cite{sadinle2019least}. MoG-HPD (Section~\ref{sec:moghpd}), which is the
$(c,\rho,\theta,\phi)=(1,1,0,0)$ base case of everything here, is essentially their construction
applied to a Mixture-of-Experts' mixture density. We do not claim it as a contribution; it is the
thing the tunable exponents are built \emph{on top of}.

\paragraph{Scaling the interval by a reported uncertainty is not ours either.} Dividing the
residual by a model-reported scale, $|y-\hat y|/\sigma(\bm x)$ (the CP-VS idea of
Section~\ref{sec:standardcp}), is standard normalized/locally-adaptive conformal prediction
\cite{lei2018distribution,papadopoulos2011regression}, and learning a quantile head directly (CQR)
is Romano et al.\ \cite{romano2019conformalized}. The one thing STA-HPD adds here is that CP-VS
hard-codes ``width grows in exact proportion to $\sigma$'' -- in the language of
Section~\ref{sec:theta}, it pins the elasticity of width in $\sigma$ to $1$. Our shape exponent $c$
simply makes that elasticity a free, tunable number, and the recurring empirical finding that the
best $c^\star$ is often below $1$ is exactly the statement that this elasticity should \emph{not}
be fixed at the CP-VS value.

\paragraph{The closest prior work, and the honest distinction.} The idea of \emph{reshaping a
conformity score after training, specifically to make the intervals narrower while keeping
coverage}, is itself already published: Boosted Conformal Prediction \cite{xie2024boosted} does
exactly this, learning a flexible correction to a base score by gradient boosting, and length-%
optimization methods \cite{kiyani2024length} target expected set size directly subject to a
validity constraint. So the \emph{goal} in this document -- post-hoc, width-minimizing, coverage-%
preserving calibration -- is not new, and it would be wrong to pitch STA-HPD/SETA-HPD as the first
to pursue it. What is different here is \emph{how narrow the family is}: instead of a boosted
ensemble or a general constrained optimizer, we use two to four interpretable exponents on a
small discrete grid, chosen so that the named calibrators a Mixture-of-Experts practitioner already
uses (Standard CP, CP-VS, MoG-HPD) sit at specific grid points (Table~\ref{tab:cheatsheet}). The
payoff is interpretability -- a fitted $(c^\star,\rho^\star,\theta^\star,\phi^\star)$ says in a
handful of numbers how the model's reported density had to be corrected -- and a mixture-density
region shape that can be genuinely multimodal (the disjoint-interval case worked through in
Section~\ref{sec:hpd}), which a boosted scalar score does not target. The cost is generality: these
methods are a low-degree-of-freedom special case, not a competitor to the full flexibility of
\cite{xie2024boosted,kiyani2024length}.

\paragraph{What is actually new, then.} Narrowly: (i) the per-component \emph{shape} exponent $c$
that powers each expert's own standard deviation before the mixture density is re-formed
(Section~\ref{sec:c}); (ii) its epistemic counterpart $\rho$ on the between-expert spread
(Section~\ref{sec:seta}); (iii) the epistemic-\emph{share} threshold field $\phi$
(Section~\ref{sec:ease}), which lets the threshold respond to how a point's uncertainty is
\emph{composed}, not only to its total; and (iv) the observation that all of Standard CP, CP-VS and
MoG-HPD are literal corners of one $(c,\rho,\theta,\phi)$ family, so a single width-driven tuning
search is guaranteed to do no worse than the best of them on its own objective
(Propositions~\ref{prop:starecovery},~\ref{prop:setarecovery},~\ref{prop:easerecovery}). The
repeated-split, mode-vote tuning of Section~\ref{sec:stafit} is a standard cross-conformal
resampling idea \cite{gibbs2021adaptive}, used here for stability rather than as a contribution in
its own right; MELD-HPD's proper-scoring-rule selector (Section~\ref{sec:meld}) is likewise a
direct application of the strictly-proper-scoring-rule framework \cite{gneiting2007strictly} to
this family's shape exponents, and its own worked comparison (Table~\ref{tab:allfive}) is included
precisely because it shows the two selection principles -- width arg-min and log-score arg-max --
can and do disagree on real data, which is the substantive content of
Remark~\ref{rem:logscoreproxy}, not merely a restatement of the framework it borrows. A careful
reader (or reviewer) will rightly ask how the numbers compare against Boosted CP directly; that
head-to-head is future work, and nothing in this document claims to beat it.

% ============================================================================
\section{Summary}
\label{sec:summary}

\subsection{Every method in this document, in one table}

\begin{table}[h]
\centering
\caption{Every calibrator introduced in this document, as a point (or sub-family) of EASE-HPD's
$(c,\rho,\theta,\phi)$ grid. ``--'' means the parameter does not exist for that method (it predates
the knob being introduced); ``n/a'' means the method does not use that grid at all, since it
selects its parameters by a different mechanism (MELD-HPD).}
\label{tab:cheatsheet}
\begin{tabular}{@{}lllll@{}}
\toprule
\textbf{Method} & $c$ & $\rho$ & $\theta$ & $\phi$ \\
\midrule
Standard CP (Sec.~\ref{sec:standardcp}) & -- & -- & -- & -- \\
CP-VS, total (Sec.~\ref{sec:standardcp}) & -- & -- & -- & -- \\
MoG-HPD (Sec.~\ref{sec:moghpd}) & $1$ & $1$ & $0$ & $0$ \\
Epistemic-scale-only HPD & $1$ & any & $0$ & $0$ \\
STA-HPD (Sec.~\ref{sec:stahpd}) & any & $1$ & any & $0$ \\
SETA-HPD (Sec.~\ref{sec:seta}) & any & any & any & $0$ \\
EASE-HPD (Sec.~\ref{sec:ease}) & any & any & any & any \\
MELD-HPD (Sec.~\ref{sec:meld}) & any\textsuperscript{$\dagger$} & any\textsuperscript{$\dagger$} &
any & n/a \\
\bottomrule
\end{tabular}
\end{table}

Standard CP and CP-VS are listed with ``--'' because they predate the level-set family entirely
(Section~\ref{sec:standardcp}); they only coincide with a grid point of $(c,\rho,\theta,\phi)$ in
the special case $K=1$ (Corollary~\ref{cor:k1closedform} and Proposition~\ref{prop:setarecovery}).
MELD-HPD's $(c,\rho)$ ($\dagger$) are drawn from the same grids $\mathcal C\times\mathcal R$ as
every other method, but selected by maximizing the log-score objective $L(c,\rho)$
(Eq.~\ref{eq:logscore}) rather than by minimizing tune-fold width -- a genuinely different
selection \emph{rule}, not a different reachable \emph{set} (Section~\ref{sec:meld}). MELD-HPD has
no $\phi$ knob at all: it is built directly on SETA-HPD's three-exponent score
(Eqs.~\ref{eq:setascore}--\ref{eq:setaregion}), not EASE-HPD's four-exponent one.

\subsection{The one-paragraph version}

A trained Mixture-of-Experts regressor reports a full Gaussian-mixture predictive density, not
just a point estimate (Section~\ref{sec:moe}). Split conformal prediction (Section~\ref{sec:cp})
turns \emph{any} fixed scoring rule built from that density into a set with an exact,
distribution-free coverage guarantee, regardless of whether the reported density is correct
(Theorem~\ref{thm:coverage}) -- so coverage is never the hard part. The hard part is width, and
the highest-density region of a density is provably the narrowest set achieving a given coverage
\emph{for that density} (Theorem~\ref{thm:hpd}), which is what MoG-HPD conformalizes
(Section~\ref{sec:moghpd}). Because the reported density is only an approximation to the truth,
STA-HPD, SETA-HPD, and EASE-HPD each add one more tunable exponent -- $c$ on the within-expert
(aleatoric) scale, $\rho$ on the between-expert (epistemic) spread, $\theta$ on how strongly the
region's threshold responds to the model's own total reported scale, and $\phi$ on how strongly it
additionally responds to that scale's \emph{composition} (Section~\ref{sec:ease}) -- each selected
by a validity-preserving, repeated-split tuning search (Section~\ref{sec:stafit}), each provably no
worse than its ancestor on the tuning objective
(Propositions~\ref{prop:starecovery},~\ref{prop:setarecovery}, and~\ref{prop:easerecovery}), and
each targeting a specific, named mismatch between the model's reported density and the true
conditional law (Section~\ref{sec:optimality}). MELD-HPD (Section~\ref{sec:meld}) departs from that
pattern in one respect: it keeps SETA-HPD's exact $(c,\rho,\theta)$ family but selects the two
density-shaping exponents $(c,\rho)$ by a proper scoring rule (maximizing calibration
log-likelihood) instead of by tune-fold width, trading the width-arg-min family's tuning-fold floor
guarantee for a different, population-level (KL-optimality) justification -- and, on real data, the
two selection principles can and do disagree (Table~\ref{tab:allfive}).

\subsection{Glossary of symbols}

\begin{center}
\begin{tabular}{@{}ll@{}}
\toprule
Symbol & Meaning \\
\midrule
$\bm x$, $y$ & input features; scalar outcome \\
$K$ & number of MoE experts \\
$\mu_k(\bm x),\sigma_k(\bm x),w_k(\bm x)$ & $k$-th expert's mean, std.\ dev., gate weight \\
$\hat y(\bm x)$ & aggregate point prediction, $\sum_k w_k\mu_k$ \\
$\sigma_a^2,\sigma_e^2,\sigma_{\mathrm{tot}}^2$ & aleatoric, epistemic, total variance \\
$\alpha$, $1-\alpha$ & target miscoverage / coverage level \\
$\mathcal D_{\mathrm{cal}},\mathcal D_{\mathrm{tune}},\mathcal D_{\mathrm{calib}}$ & calibration
set and its tune/calib split \\
$S_i$, $\hat q$ / $\hat t$ & nonconformity score; conformal threshold \\
$f_{\mathrm{mix}}$ & the model's raw Gaussian-mixture density \\
$R_\alpha(f)$ & highest-density region of $f$ at level $\alpha$ \\
$G$ & calibration-set geometric-mean expert scale (Eq.~\ref{eq:Gdef}) \\
$c$ & aleatoric shape exponent (Eq.~\ref{eq:sigmashape}) \\
$\theta$ & threshold-field exponent (Eq.~\ref{eq:stascore}) \\
$\rho$ & epistemic scale exponent (Eq.~\ref{eq:etamu}) \\
$\tilde\sigma_k(\bm x;c)$, $\tilde\mu_k(\bm x;\rho)$ & reshaped expert scale / shifted expert mean \\
$A(\bm x;c)$, $E(\bm x;\rho)$ & aleatoric, epistemic terms of $\sigma_{\mathrm{tot}}^2$ (Eq.~\ref{eq:aeterms}) \\
$s(\bm x;c,\rho)$, $\bar s(c,\rho)$ & epistemic share (Eq.~\ref{eq:share}); its calibration-set mean \\
$\phi$ & epistemic-share threshold-field exponent (Eq.~\ref{eq:easescore}) \\
$L(c,\rho)$ & mean tune-fold log-density objective (Eq.~\ref{eq:logscore}) \\
$\hat\rho$ & MELD-HPD's closed-form epistemic-reliability diagnostic (Eq.~\ref{eq:closedrho}) \\
\bottomrule
\end{tabular}
\end{center}

% ============================================================================
\begin{thebibliography}{12}

\bibitem{vovk2005algorithmic}
V.~Vovk, A.~Gammerman, and G.~Shafer, \emph{Algorithmic Learning in a Random World}. Springer,
2005.

\bibitem{lei2018distribution}
J.~Lei, M.~G'Sell, A.~Rinaldo, R.~J.~Tibshirani, and L.~Wasserman, ``Distribution-free predictive
inference for regression,'' \emph{J. Amer. Statist. Assoc.}, vol.~113, no.~523, pp.~1094--1111,
2018.

\bibitem{romano2019conformalized}
Y.~Romano, E.~Patterson, and E.~Cand\`es, ``Conformalized quantile regression,'' in \emph{Advances
in Neural Information Processing Systems (NeurIPS)}, 2019.

\bibitem{izbicki2020cd}
R.~Izbicki, G.~Shimizu, and R.~B.~Stern, ``Flexible distribution-free conditional predictive bands
using density estimators,'' in \emph{Proc. International Conference on Artificial Intelligence and
Statistics (AISTATS)}, 2020.

\bibitem{sadinle2019least}
M.~Sadinle, J.~Lei, and L.~Wasserman, ``Least ambiguous set-valued classifiers with bounded error
levels,'' \emph{J. Amer. Statist. Assoc.}, vol.~114, no.~525, pp.~223--234, 2019.

\bibitem{boxtiao1973bayesian}
G.~E.~P.~Box and G.~C.~Tiao, \emph{Bayesian Inference in Statistical Analysis}. Addison-Wesley,
1973.

\bibitem{hyndman1996computing}
R.~J.~Hyndman, ``Computing and graphing highest density regions,'' \emph{The American
Statistician}, vol.~50, no.~2, pp.~120--126, 1996.

\bibitem{jordan1994hierarchical}
M.~I.~Jordan and R.~A.~Jacobs, ``Hierarchical mixtures of experts and the EM algorithm,''
\emph{Neural Computation}, vol.~6, no.~2, pp.~181--214, 1994.

\bibitem{gibbs2021adaptive}
I.~Gibbs and E.~Cand\`es, ``Adaptive conformal inference under distribution shift,'' in
\emph{Advances in Neural Information Processing Systems (NeurIPS)}, 2021.

\bibitem{gneiting2007strictly}
T.~Gneiting and A.~E.~Raftery, ``Strictly proper scoring rules, prediction, and estimation,''
\emph{J. Amer. Statist. Assoc.}, vol.~102, no.~477, pp.~359--378, 2007.

\bibitem{gneiting2007strictly}
T.~Gneiting and A.~E.~Raftery, ``Strictly proper scoring rules, prediction, and estimation,''
\emph{J. Amer. Statist. Assoc.}, vol.~102, no.~477, pp.~359--378, 2007.

\bibitem{papadopoulos2011regression}
H.~Papadopoulos, V.~Vovk, and A.~Gammerman, ``Regression conformal prediction with nearest
neighbours,'' \emph{J. Artificial Intelligence Research}, vol.~40, pp.~815--840, 2011.

\bibitem{xie2024boosted}
R.~Xie, R.~Foygel~Barber, and E.~J.~Cand\`es, ``Boosted conformal prediction intervals,'' in
\emph{Advances in Neural Information Processing Systems (NeurIPS)}, 2024.

\bibitem{kiyani2024length}
S.~Kiyani, G.~Pappas, and H.~Hassani, ``Length optimization in conformal prediction,'' in
\emph{Advances in Neural Information Processing Systems (NeurIPS)}, 2024.

\end{thebibliography}

\end{document}
